\chapter{x86-64 汇编语法总览}

\section{问题与目标}

前面几章已经建立了三条基础线：

\begin{lstlisting}
C++ source -> compiler -> assembly text -> machine code bytes
object/type/layout -> address formula -> memory operand
instruction bytes -> CPU state transition -> performance behavior
\end{lstlisting}

现在正式进入 x86-64 汇编。初学阶段常见两种相反误区：一种把汇编当成“更难看的 C”，另一种把汇编当成不可读的符号堆。两者都不准确。汇编不是 C，它更接近 ISA 暴露给软件的机器动作；但汇编也不是随机符号，只要掌握语法、寄存器、操作数、指令宽度和 ABI 上下文，就可以逐条推导。

本章目标不是一次背完所有 x86 指令，而是建立读汇编的基本方法。本章收束时，应能做到：

\begin{itemize}
  \item 看懂一条 Intel 语法汇编指令的基本结构。
  \item 区分立即数、寄存器操作数、内存操作数。
  \item 识别通用寄存器及其 64/32/16/8 位别名。
  \item 理解 \code{byte ptr}、\code{dword ptr}、\code{qword ptr} 等宽度标记。
  \item 从 \code{[base + index*scale + displacement]} 还原 C++ 数组和结构体访问。
  \item 用“读哪些状态、写哪些状态”的方式注释汇编。
  \item 比较 \code{-O0} 和 \code{-O3} 汇编为什么差异巨大。
  \item 知道 Intel 语法和 AT\&T 语法的关键差异，避免读资料时混淆。
\end{itemize}

\begin{keyidea}
读汇编的核心不是把每条指令机械翻译成一行 C++，而是恢复数据流、控制流和地址公式。C++ 源码只是汇编的一种可能来源；真正可靠的依据，是逐条说明指令读写了哪些寄存器、内存和 flags。
\end{keyidea}

\section{读汇编前先分清四个层次}

汇编入门容易混乱，不是因为语法本身多难，而是因为几个层次的知识经常同时出现。看到一行汇编时，里面至少包含四类信息：汇编语法、ISA 语义、ABI 约定和编译器选择。把这四类信息分清，后面读任何函数都会更稳定。

第一类是汇编语法。它回答“这行文本怎么读”。例如 Intel 语法中通常写成目的操作数在前、源操作数在后；方括号表示内存操作数；\code{dword ptr} 说明内存访问宽度。语法本身不解释为什么第一个参数在 \register{rdi}，也不解释一条指令在某个 CPU 上快不快。

第二类是 ISA 语义。它回答“这条指令对架构状态做什么”。例如 \instruction{add eax, esi} 读取 \register{eax} 和 \register{esi}，把低 32 位加法结果写回 \register{eax}，并更新一部分 flags；\instruction{ret} 从栈顶取返回地址，改变 \register{rsp} 和 \register{rip}。ISA 语义关心软件可见状态，不要求你知道 CPU 内部把它拆成几个微操作。

第三类是 ABI 约定。它回答“函数之间如何协作”。为什么第一个整数参数常在 \register{edi}，为什么返回值在 \register{eax}，为什么某些寄存器由调用者保存、某些由被调用者保存，这不是 \instruction{mov} 或 \instruction{add} 本身决定的，而是 System V AMD64 ABI 决定的。同一条指令在不同函数位置有不同含义，常常就是因为 ABI 上下文不同。

第四类是编译器选择。它回答“为什么这段源码生成了这种形状”。编译器可以选择内联、删除、展开、合并、使用 \instruction{lea} 代替乘法、使用条件移动代替分支，也可以因为调试构建而保留大量栈槽。编译器选择受优化级别、目标 CPU、语言规则、别名信息和可见函数体影响。你不能只凭一段汇编判断源码原来一定长什么样，只能判断这段汇编实际做了什么。

\begin{mentalmodel}
看到一行汇编时，按顺序问：语法上谁是目的操作数，ISA 上它读写什么状态，ABI 上这些寄存器在函数入口或出口代表什么，编译器为什么可能选择这种写法。四个问题分开回答，错误会少很多。
\end{mentalmodel}

\section{汇编不是更低级的 C}

很多初学者会把汇编读成“一条指令对应一行 C++”。这在极小例子里偶尔能用，但作为方法是危险的。C++ 是带类型、作用域、对象生命周期、别名规则和未定义行为的语言；汇编面对的是寄存器、内存地址、flags 和控制流。两者之间经过优化器，形状可以完全不同。

一个 C++ 局部变量未必对应内存中的一个槽。它可能一直在寄存器里，也可能被常量传播后消失，也可能和另一个临时值复用同一个寄存器。相反，\code{-O0} 汇编里的许多栈槽也不代表真实性能路径中的对象，只是为了调试器能在源代码行之间看到变量。把调试构建的栈槽当成语言对象的真实运行形态，会误导后面的性能判断。

一个 C++ 表达式也未必对应一条指令。表达式可以被拆成多条指令，也可以和相邻表达式合并，还可以被代数化简。比如乘以常数可能用 \instruction{lea}，比较可能只留下 flags，条件表达式可能变成 \instruction{cmov}，数组访问可能和加法合并成带内存操作数的指令。读汇编时不要急着把每条指令翻译回某一行源码，而要先恢复数据如何流动、地址如何形成、控制流如何选择。

最可靠的读法是把汇编看成状态转移。每条指令都在读某些状态、写某些状态。状态可以是通用寄存器，可以是内存，可以是 flags，也可以是 \register{rip} 和 \register{rsp}。如果你能完整说明这些读写，再结合函数入口的 ABI 约定，就能解释这段机器码。至于它最初来自哪一行 C++，要在最后再谨慎推断。

\begin{warningbox}
不要把“恢复 C-like 伪代码”当作第一步。伪代码是分析结果，不是分析依据。先做寄存器含义、内存地址、访问宽度和 flags 依赖，再写伪代码。
\end{warningbox}

\section{先生成一份你能控制的汇编}

学习汇编不能只看网上片段。你要从自己写的 C++ 生成汇编，知道编译命令、优化级别、目标语法和反汇编工具。先创建 \filepath{labs/ch10_assembly_syntax/minimal.cpp}：

\begin{lstlisting}[language=C++]
extern "C" int add3(int x, int y, int z)
{
    return x + y + z;
}

extern "C" int load_i32(int const* p, int i)
{
    return p[i];
}
\end{lstlisting}

生成 Intel 语法汇编：

\begin{lstlisting}[language=bash]
clang++ -std=c++20 -O3 -S -masm=intel minimal.cpp -o minimal.s
\end{lstlisting}

生成目标文件并反汇编：

\begin{lstlisting}[language=bash]
clang++ -std=c++20 -O3 -c minimal.cpp -o minimal.o
objdump -d -Mintel minimal.o > minimal.objdump.txt
\end{lstlisting}

如果你的系统主要使用 LLVM 工具链，也可以用：

\begin{lstlisting}[language=bash]
llvm-objdump -d --x86-asm-syntax=intel minimal.o > minimal.objdump.txt
\end{lstlisting}

\code{minimal.s} 是汇编器输入，通常包含标签、伪指令、调试/展开信息等。\code{minimal.objdump.txt} 是从目标文件机器码反汇编出来的文本，会显示机器码字节和指令地址。学习时两者都要看：

\begin{itemize}
  \item \code{.s} 更接近编译器输出，适合看函数边界和汇编器伪指令。
  \item \code{objdump} 更接近最终二进制，适合看机器码字节、真实指令地址和链接前符号。
\end{itemize}

这两份文件回答的问题不同。\code{.s} 适合观察编译器想让汇编器看到什么，包括函数标签、可见性、段切换、对齐和一些调试或展开信息。反汇编适合观察目标文件里真实存在的指令字节，尤其适合验证指令长度、相对跳转、重定位位置和符号地址。高质量报告最好同时保留两者，因为它们能互相校验。

还要保存完整编译命令。没有命令，汇编没有可复现性。优化级别、目标 CPU、是否启用位置无关代码、是否保留帧指针、是否启用链接时优化、使用 GCC 还是 Clang，都会改变输出。一个报告只贴几行汇编，却不说明命令，读者无法判断这些指令是调试路径、发布路径，还是某个特定工具链版本的偶然结果。

生成汇编时也要避免一次引入太多库和模板。第一份练习最好使用 \code{extern "C"} 的小函数，避免 C++ 名字改编、异常表、标准库初始化和模板实例化把观察对象淹没。先控制变量，再扩大复杂度：从小函数开始，逐步加入结构体、数组、循环、函数调用和标准库代码。

\begin{warningbox}
不同编译器、版本、优化级别、目标 CPU、命令行参数都会生成不同汇编。本书的汇编是“合理示例”，不是要求你的机器逐字一致。学习重点是解释自己的输出。
\end{warningbox}

\section{汇编器到底做什么}

很多读者第一次看到 \filepath{.s} 文件，会把它理解成“CPU 直接执行的文本”。这不准确。汇编器的任务是把汇编文本翻译成目标文件中的机器码、符号和重定位信息。CPU 最终执行的是内存中的机器码字节，不是 \filepath{.s} 文件。

从编译器到 CPU，中间至少有三个可观察产物：

\begin{center}
\begin{tabular}{p{0.20\linewidth}p{0.32\linewidth}p{0.34\linewidth}}
\toprule
产物 & 谁产生或消费 & 主要用途 \\
\midrule
\filepath{.s} & 编译器产生，汇编器消费 & 人类可读的汇编文本、标签、伪指令、段信息 \\
\filepath{.o} & 汇编器产生，链接器消费 & 机器码、符号表、重定位、调试/展开信息 \\
可执行文件 & 链接器产生，加载器消费 & 最终程序映像、动态链接信息、入口和段映射 \\
\bottomrule
\end{tabular}
\end{center}

汇编器做的事包括：

\begin{itemize}
  \item 把助记符和操作数编码成机器码字节。
  \item 计算同一汇编文件内标签之间的相对位移。
  \item 把函数名、全局对象名等记录为符号。
  \item 对当前还不能确定的地址生成重定位记录。
  \item 按伪指令切换 section、设置对齐、生成常量数据和元数据。
\end{itemize}

这解释了为什么 \filepath{.s} 文件里会有很多不是 CPU 指令的行。比如 \code{.text}、\code{.globl}、\code{.type}、\code{.size}、\code{.cfi_startproc}、\code{.p2align} 都是给汇编器、链接器、调试器或异常展开机制看的元信息。它们会影响目标文件，但不是 CPU 在函数热路径中逐条执行的指令。

例如：

\begin{lstlisting}
.text
.globl add2
.type add2,@function
add2:
    mov eax, edi
    add eax, esi
    ret
.size add2, .-add2
\end{lstlisting}

这里真正的指令只有三条：\instruction{mov}、\instruction{add}、\instruction{ret}。标签 \code{add2:} 是地址名字，不是指令；\code{.globl} 让符号对链接器可见；\code{.type} 和 \code{.size} 帮助工具知道这是函数以及函数范围。

\begin{keyidea}
读汇编文件时要先把“CPU 指令”和“汇编器/链接器元信息”分开。伪指令不是没用，但它回答的是目标文件组织和工具链问题，不是每周期执行问题。
\end{keyidea}

\section{section：代码、常量和未初始化数据怎样分区}

汇编文件不是一串扁平指令。汇编器输出目标文件时，会把不同性质的内容放入不同 section。常见 section 包括 \code{.text}、\code{.rodata}、\code{.data} 和 \code{.bss}。这些名字看起来像格式细节，其实会影响链接、加载、权限和运行时内存映像。

\begin{center}
\begin{tabular}{p{0.22\linewidth}p{0.31\linewidth}p{0.32\linewidth}}
\toprule
section & 常见内容 & 运行时性质 \\
\midrule
\code{.text} & 函数机器码、跳转填充、只执行的代码 & 通常可读可执行，不应可写 \\
\code{.rodata} & 字符串字面量、只读常量表、跳转表 & 通常可读不可写 \\
\code{.data} & 有初始值的全局变量和静态变量 & 可读可写，文件中保存初始值 \\
\code{.bss} & 零初始化的全局变量和静态变量 & 可读可写，文件中通常不保存大块零字节 \\
\bottomrule
\end{tabular}
\end{center}

例如：

\begin{lstlisting}
.text
.globl read_answer
.type read_answer,@function
read_answer:
    mov eax, dword ptr [rip + answer]
    ret

.data
.globl answer
answer:
    .long 42

.section .rodata
message:
    .asciz "hello"
\end{lstlisting}

这里 \code{read_answer} 的指令在 \code{.text}；\code{answer} 的初始值在 \code{.data}；\code{message} 的字节在只读数据 section。CPU 执行时不会“执行 \code{.data}”。但链接器和加载器会根据 section 属性把它们映射到进程虚拟地址空间的不同区域，并设置合适的页权限。

这也解释了为什么同样是字节，工程含义完全不同。\code{.text} 里的字节会被反汇编器当作指令候选；\code{.rodata} 里的字节可能是字符串、浮点常量、查表数据或跳转表；\code{.bss} 里的对象在目标文件中可能只记录大小，不真的存储一堆零。

读目标文件时，如果不知道当前字节属于哪个 section，就很容易把数据当代码、把常量当指令，或者把未初始化存储误认为文件里漏了内容。

\subsection{section 切换不是运行时跳转}

伪指令 \code{.text}、\code{.data}、\code{.section .rodata} 的作用是告诉汇编器“接下来输出到哪个区域”。它们不是运行时控制流。下面这段文本：

\begin{lstlisting}
.text
f:
    ret

.data
x:
    .long 7
\end{lstlisting}

不是说函数 \code{f} 返回后会跳到 \code{x}，也不是说 CPU 会执行 \code{.data}。汇编器只是在构造目标文件时把 \code{ret} 编码到代码 section，把 4 字节整数常量放到数据 section。运行时控制流由 \instruction{call}、\instruction{ret}、\instruction{jmp}、条件跳转和异常/信号等机制决定，不由汇编文本中 section 出现的先后决定。

这个区分对阅读编译器输出尤其重要。一个 \filepath{.s} 文件可能先写某个函数，随后切到 \code{.rodata} 输出字符串，再切回 \code{.text} 输出下一个函数。文本顺序只是汇编器输入顺序，最终链接器还可能重排 section、合并相同常量、丢弃未引用 section、插入对齐填充或把热冷代码分开。报告中如果要说明运行时地址顺序，应使用链接后反汇编或 section header，而不是只凭 \filepath{.s} 的文本顺序。

\subsection{数据伪指令写的是字节，不是高级类型}

汇编里的数据伪指令常见有 \code{.byte}、\code{.word}、\code{.long}、\code{.quad}、\code{.ascii}、\code{.asciz}、\code{.zero}。它们告诉汇编器输出特定宽度的字节序列：

\begin{lstlisting}
.byte  1
.word  0x0203
.long  0x04050607
.quad  0x08090a0b0c0d0e0f
.asciz "ok"
.zero  16
\end{lstlisting}

这些伪指令不携带完整 C++ 类型语义。\code{.long 1} 只是输出 4 字节整数编码，可能表示 \code{int}，也可能表示无符号整数、位掩码、浮点位模式的一部分、跳转表项或某个结构体字段。是否能把它解释成 C++ 类型，要看 section、符号、对齐、引用它的指令、调试信息和源码上下文。

字符串也一样。\code{.asciz} 通常输出带结尾零字节的 C 风格字符串；\code{.ascii} 不自动追加结尾零。看到 \code{.asciz "hello"}，可以推断它很可能服务于字符串字面量或 C 接口，但仍要看是否有代码用 \instruction{lea rdi, [rip + ...]} 取它的地址，或者是否被作为常量表的一部分引用。

\begin{warningbox}
不要把数据伪指令直接当作“变量声明”。汇编层只输出字节、对齐和符号；高级语言的类型、生命周期、const 语义和所有权，要靠源码和 ABI 共同解释。
\end{warningbox}

\section{对齐和填充：为什么会出现看似多余的字节}

目标文件中经常出现对齐伪指令，例如 \code{.p2align}、\code{.align}、\code{.balign}。编译器可能在函数入口、循环入口、常量表或对象之间插入填充，让后续地址满足某种边界要求。

\begin{lstlisting}
.p2align 4
hot_loop:
    add eax, dword ptr [rdi + rcx*4]
    inc rcx
    cmp rcx, rsi
    jl hot_loop
\end{lstlisting}

\code{.p2align 4} 常表示按 \(2^4 = 16\) 字节边界对齐。汇编器可能在 \code{hot_loop} 前插入若干字节。若填充位于 \code{.text}，它可能使用一种或多种 \instruction{nop} 编码；若填充位于数据 section，可能是零字节或指定填充值。

对齐的目的有几类。代码入口对齐可能改善取指和分支目标布局；循环对齐可能减少跨取指块边界；数据对齐可能满足类型要求或让后续 load/store 更自然；向量数据可能需要更严格的边界。第一册不深入 CPU 前端和 cache line 性能，但要先知道：填充字节通常是工具链有意放置的布局内容，不一定是“无用垃圾”。

反汇编时，对齐填充会带来两个常见误判。第一，把函数之间的 \instruction{nop} 当作源代码逻辑。普通 \instruction{nop} 不改变程序语义，它可能只是对齐填充。第二，把数据 section 中的填充字节反汇编成指令。反汇编器如果从错误位置开始解释字节流，任何数据都可能被显示成看似合法的 x86 指令。x86 是变长指令集，很多随机字节序列都能被解释出某种指令文本，所以“反汇编器显示了助记符”本身不证明这些字节真会被执行。

\begin{deepdive}
后续学习性能时，对齐不应被神化。对齐可能影响取指、解码、cache line、分页和向量访问，但具体效果依赖代码位置、循环形状、CPU、数据规模和运行路径。当前阶段只要求你能识别对齐伪指令和填充字节，并在报告里把它们与真实业务指令分开。
\end{deepdive}

\section{局部标签、全局符号和可见性}

汇编文件中有些名字只服务于当前文件内部跳转，有些名字要暴露给链接器和其他目标文件。前者通常是局部标签，后者通常是全局符号。分清二者，可以避免把编译器内部控制流标签误认为公共 API。

编译器生成的局部标签常见形式包括 \code{.LBB0_1}、\code{.Ltmp3}、\code{.L.str}。不同工具链命名不同，但一般以 \code{.L} 开头的符号默认不会作为普通全局符号导出。它们可能标记基本块、常量池、异常处理范围、调试临时位置或跳转表位置。

\begin{lstlisting}
.globl sum_positive
.type sum_positive,@function
sum_positive:
    xor eax, eax
    test esi, esi
    jle .Ldone
.Lloop:
    add eax, dword ptr [rdi]
    add rdi, 4
    dec esi
    jne .Lloop
.Ldone:
    ret
\end{lstlisting}

这里 \code{sum_positive} 是外部可见函数符号；\code{.Lloop} 和 \code{.Ldone} 是函数内部控制流标签。机器码中的 \instruction{jne .Lloop} 通常会变成相对位移，运行时不需要保存字符串 \code{.Lloop}。如果链接器最终丢弃局部符号，程序仍能运行，因为跳转目标已经被编码成位移或重定位修补后的地址。

\subsection{\code{.globl} 只解决链接可见性}

\code{.globl name} 的意思是让 \code{name} 成为链接器可见的全局符号。它不自动保证这个符号是函数，不自动保证 ABI 正确，也不自动保证 C++ 端能按期望名字找到它。函数符号通常还配合 \code{.type name,@function} 和 \code{.size name, .-name}，让工具知道符号类型和范围。

对手写汇编来说，只写标签而忘记 \code{.globl}，C++ 代码可能链接不到该函数；写了 \code{.globl} 但 C++ 声明没有 \code{extern "C"}，也可能因为 C++ 名字修饰导致符号名不匹配；符号名匹配但参数寄存器、栈对齐或 callee-saved 寄存器规则错了，程序仍然可能崩溃。符号可见性只是跨文件调用的第一道门，ABI 正确性还要在后面章节系统检查。

\subsection{符号范围影响工具输出}

使用 \code{nm}、\code{readelf -s} 或 \code{objdump -t} 时，你会看到符号的绑定、类型、section 和地址偏移。一个典型目标文件可能同时包含：

\begin{lstlisting}
0000000000000000 T sum_positive
                 U printf
0000000000000000 r .L.str
\end{lstlisting}

这类输出中，\code{T} 常表示定义在代码 section 的全局文本符号，\code{U} 表示 undefined symbol，需要链接器从别处解析，\code{r} 可能表示只读数据中的局部符号。不同工具输出格式会变，但核心问题一样：这个名字是本文件定义，还是外部引用；是函数、对象，还是 section 内局部位置；链接后是否仍然可见。

读符号表时，不要只看名字。还要看绑定属性、类型、所在 section、大小和是否 undefined。一个 \code{printf} 出现在符号表里，通常说明当前目标文件引用了外部 \code{printf}，并不说明 \code{printf} 的代码已经在当前目标文件中。一个 \code{.L.str} 出现在只读 section，通常说明它是字符串或常量位置，不应被当作可调用函数。

\section{标签、符号和地址不是一回事}

汇编里经常看到名字，例如函数标签、局部标签、全局对象名和跳转目标。它们都可能显示成“符号”，但含义要分清。

标签是汇编文本中的一个位置标记：

\begin{lstlisting}
loop:
    add eax, dword ptr [rdi + rcx*4]
    inc rcx
    cmp rcx, rsi
    jl loop
\end{lstlisting}

\code{loop} 标记某条指令的地址。条件跳转 \instruction{jl loop} 在机器码里通常不会保存字符串 \code{loop}，而是保存从当前指令到目标位置的相对位移。汇编器负责把标签名转成位移或重定位。

符号是目标文件层面记录的名字。函数名、全局变量名、外部引用都可能出现在符号表中。使用 \code{nm -C} 或 \code{readelf -s} 可以看到符号。一个标签不一定成为外部可见符号；一个符号也可能没有完整函数体，例如 undefined symbol 只表示当前目标文件引用了外部定义。

地址是运行时或目标文件中的位置数值。目标文件中的地址常常只是 section 内偏移；最终可执行文件里链接器会分配虚拟地址；启用 PIE 和 ASLR 后，运行时装载基址还可能变化。报告里说“函数在某地址”，必须说明这是目标文件偏移、可执行文件虚拟地址，还是当前进程运行时地址。

这个区分对调试非常重要。GDB 中的 \code{break add2} 是按符号下断点；\code{break *0x401140} 是按地址下断点；汇编中的 \code{add2:} 是标签；反汇编里的 \code{0000000000000000 <add2>:} 是工具把地址和符号名结合显示出来。它们有关联，但不是同一个概念。

\begin{warningbox}
不要把“看见名字”理解成“机器码里保存了这个名字”。大多数普通指令执行时只处理位模式、寄存器编号、立即数、位移和地址。名字主要服务于汇编、链接、调试和报告。
\end{warningbox}

\section{一条汇编指令的基本结构}

Intel 语法里，一条普通指令大致长这样：

\begin{lstlisting}
mnemonic destination, source
\end{lstlisting}

例如：

\begin{lstlisting}
mov eax, edi
add eax, esi
ret
\end{lstlisting}

这段代码可以来自：

\begin{lstlisting}[language=C++]
extern "C" int add2(int x, int y)
{
    return x + y;
}
\end{lstlisting}

在 System V AMD64 ABI 下，前两个 \code{int} 参数在 \register{edi} 和 \register{esi}，返回值在 \register{eax}。逐条解释：

\begin{center}
\begin{tabular}{p{0.21\linewidth}p{0.33\linewidth}p{0.34\linewidth}}
\toprule
指令 & 读什么 & 写什么 \\
\midrule
\code{mov eax, edi} & \register{edi} & \register{eax}，同时清零 \register{rax} 高 32 位 \\
\code{add eax, esi} & \register{eax}、\register{esi} & \register{eax} 和部分 flags \\
\code{ret} & 栈顶返回地址、\register{rsp} & \register{rip}、\register{rsp} \\
\bottomrule
\end{tabular}
\end{center}

这里有一个初学阶段必须接受的事实：\instruction{add eax, esi} 的目的操作数既是输入也是输出。它不是“产生一个新变量”，而是把结果写回 \register{eax}：

\begin{lstlisting}
eax = eax + esi
\end{lstlisting}

因此读汇编时，你应该维护一个“当前每个寄存器代表什么”的笔记。比如：

\begin{lstlisting}
before:
    edi = x
    esi = y

mov eax, edi
    eax = x

add eax, esi
    eax = x + y

ret
    return eax
\end{lstlisting}

这个例子虽然小，但已经包含读汇编的核心动作。函数入口时，\register{edi} 和 \register{esi} 的含义不是从指令本身来的，而是从 ABI 和函数签名来的。第一条 \instruction{mov} 之后，\register{eax} 才开始代表 \code{x}。第二条 \instruction{add} 之后，\register{eax} 的含义变成 \code{x + y}。最后 \instruction{ret} 并不“返回 \register{eax}”这个动作本身；返回值放在 \register{eax} 是 ABI 约定，\instruction{ret} 做的是控制流返回。

因此，逐条注释时不能只写“移动”“相加”“返回”。这种注释没有信息量。好的注释要写出入口假设、读写状态和状态含义的变化。尤其要注意目的操作数既可能是输入也是输出，例如 \instruction{add}、\instruction{sub}、\instruction{and}、\instruction{or} 都会读取旧目的值，再写回新目的值。如果忘记这一点，就会把数据依赖链看断。

\begin{deepdive}
现代性能分析非常重视依赖链。即使两条指令看起来都很短，只要后一条必须等待前一条写出结果，就不能完全并行。读汇编时维护寄存器含义，不只是为了理解语义，也是为了后续判断 latency、throughput 和乱序执行空间打基础。
\end{deepdive}

\section{操作数：立即数、寄存器和内存}

x86-64 指令的操作数常见三类：立即数、寄存器、内存。

\subsection{立即数}

\term{immediate}{立即数} 是编码在指令里的常量：

\begin{lstlisting}
mov eax, 42
add rdi, 16
and eax, 255
\end{lstlisting}

立即数不是内存地址，也不需要额外 load。它是机器码字节的一部分。比如 \code{add rdi, 16} 中的 \code{16} 会被编码进指令。

\subsection{寄存器操作数}

寄存器操作数直接读写 CPU 的架构寄存器：

\begin{lstlisting}
mov rax, rdi
add eax, esi
xor ecx, ecx
\end{lstlisting}

寄存器访问通常比内存访问快得多，但寄存器数量有限。编译器优化很大一部分工作，就是尽量把热数据保存在寄存器中，减少 load/store。

\subsection{内存操作数}

方括号表示内存操作数：

\begin{lstlisting}
mov eax, dword ptr [rdi]
mov eax, dword ptr [rdi + rsi*4]
mov qword ptr [rsp + 8], rax
\end{lstlisting}

方括号里的表达式计算地址；\code{byte ptr}、\code{dword ptr}、\code{qword ptr} 表示访问宽度。注意：

\begin{lstlisting}
mov eax, edi        ; register -> register
mov eax, [rdi]      ; memory at address rdi -> register
\end{lstlisting}

少一个方括号，语义完全不同。第一条把 \register{edi} 的值复制到 \register{eax}；第二条把 \register{rdi} 当作地址，从该地址读取内存。

\subsection{一般不能内存到内存}

多数普通 x86 指令不能同时有两个内存操作数。例如下面这种通常不合法：

\begin{lstlisting}
mov dword ptr [rdi], dword ptr [rsi]    ; invalid for ordinary mov
\end{lstlisting}

需要用寄存器中转：

\begin{lstlisting}
mov eax, dword ptr [rsi]
mov dword ptr [rdi], eax
\end{lstlisting}

这对性能分析很重要。源码里看起来像一次赋值，机器层面可能是一次 load 加一次 store。

\subsection{内存操作数不是对象本身}

方括号里的表达式只计算地址，访问的对象含义要靠上下文恢复。看到 \code{[rdi + rsi*4]}，只能先说它访问了从 \register{rdi} 加 \register{rsi} 乘 4 得到的地址，访问宽度由指令或宽度标记决定。它可能是 \code{int} 数组元素，也可能是某个结构体数组里的字段，也可能是编译器临时布置的数据表。只有结合函数签名、类型布局、访问宽度、相邻指令和源码，才能进一步推断它对应什么对象。

这个原则能避免许多错误。初学者看到乘 4 就直接说“这是 \code{int} 数组”，但乘 4 也可能来自四字节字段、索引表、浮点数组或某种压缩布局。看到 \code{+ 8} 也不能马上说“这是第三个元素”，它可能是结构体字段偏移、栈槽偏移、对象头之后的数据，或调用约定要求的保存位置。汇编里的地址公式是证据，C++ 对象解释是推论，二者不要混为一谈。

另一个常见错误是把内存操作数都当成“慢”。内存访问的成本取决于是否命中 cache、地址是否已知、是否和前后指令存在依赖、是否与 store 冲突、是否跨 cache line 或页面等。第 6 章不展开这些微架构细节，但你要先养成准确标注 load/store 的习惯。只有先知道哪里真正访问内存，后续才谈得上判断内存层级成本。

\section{通用寄存器和别名}

x86-64 有 16 个常用通用整数寄存器。每个寄存器可以用不同宽度的名字访问低位部分：

\begin{center}
\begin{tabular}{lllll}
\toprule
64 位 & 32 位 & 16 位 & 8 位低位 & 常见 ABI 含义 \\
\midrule
\register{rax} & \register{eax} & \register{ax} & \register{al} & 返回值、临时 \\
\register{rbx} & \register{ebx} & \register{bx} & \register{bl} & callee-saved \\
\register{rcx} & \register{ecx} & \register{cx} & \register{cl} & 第 4 参数、移位计数 \\
\register{rdx} & \register{edx} & \register{dx} & \register{dl} & 第 3 参数、乘除辅助 \\
\register{rsi} & \register{esi} & \register{si} & \register{sil} & 第 2 参数 \\
\register{rdi} & \register{edi} & \register{di} & \register{dil} & 第 1 参数 \\
\register{rbp} & \register{ebp} & \register{bp} & \register{bpl} & 帧指针或 callee-saved \\
\register{rsp} & \register{esp} & \register{sp} & \register{spl} & 栈指针 \\
\bottomrule
\end{tabular}
\end{center}

\register{r8} 到 \register{r15} 也有类似名字：

\begin{center}
\begin{tabular}{llll}
\toprule
64 位 & 32 位 & 16 位 & 8 位 \\
\midrule
\register{r8} & \register{r8d} & \register{r8w} & \register{r8b} \\
\register{r9} & \register{r9d} & \register{r9w} & \register{r9b} \\
\register{r10} & \register{r10d} & \register{r10w} & \register{r10b} \\
\register{r11} & \register{r11d} & \register{r11w} & \register{r11b} \\
\register{r12} & \register{r12d} & \register{r12w} & \register{r12b} \\
\register{r13} & \register{r13d} & \register{r13w} & \register{r13b} \\
\register{r14} & \register{r14d} & \register{r14w} & \register{r14b} \\
\register{r15} & \register{r15d} & \register{r15w} & \register{r15b} \\
\bottomrule
\end{tabular}
\end{center}

\subsection{32 位写入会清零高 32 位}

x86-64 一个非常重要的规则是：写 32 位寄存器会把对应 64 位寄存器高 32 位清零。

\begin{lstlisting}
mov eax, edi
\end{lstlisting}

这条指令写 \register{eax}，也会让 \register{rax} 的高 32 位变成 0。可以这样理解：

\begin{lstlisting}
rax = zero_extend_64(edi)
\end{lstlisting}

但写 16 位或 8 位寄存器不会自动清零高位：

\begin{lstlisting}
mov ax, di      ; only low 16 bits are written
mov al, dil     ; only low 8 bits are written
\end{lstlisting}

因此部分寄存器写入可能保留旧值的高位。现代编译器通常更喜欢用 32 位写入来打破不必要的旧值依赖。

\begin{warningbox}
不要把 \register{eax}、\register{ax}、\register{al} 当成完全独立寄存器。它们是 \register{rax} 的不同宽度视图。读汇编时必须知道一条指令到底写了多少位。
\end{warningbox}

寄存器别名还会影响你对“值是否仍然存在”的判断。写 \register{eax} 会清零 \register{rax} 高 32 位，所以之后读 \register{rax} 得到的是零扩展后的 64 位值。写 \register{al} 只改变最低 8 位，高位仍然保留旧内容。假如你在笔记里把 \register{al} 写成独立变量，就会忘记 \register{rax} 的其余位仍然有历史值。

编译器倾向于使用 32 位写入清零寄存器，也和依赖有关。写完整的 32 位子寄存器能让硬件知道高 32 位确定为零，减少对旧 \register{rax} 的依赖。某些老旧微架构上，部分寄存器写入和随后读取完整寄存器可能引发额外依赖或合并成本。现代编译器和 CPU 已经处理了很多历史问题，但作为阅读者，你仍然要知道：宽度不是排版细节，它影响语义和依赖。

\begin{mentalmodel}
寄存器不是变量名，而是一组可用不同宽度观察和写入的位。读汇编时，先确定写入宽度，再决定整个 64 位寄存器的新含义。
\end{mentalmodel}

\section{操作数宽度：机器到底读写几个字节}

Intel 语法中，访问内存时常用宽度标记：

\begin{center}
\begin{tabular}{lll}
\toprule
标记 & 字节数 & 常见 C++ 类型或用途 \\
\midrule
\code{byte ptr} & 1 & \code{char}、\code{std::uint8_t} \\
\code{word ptr} & 2 & \code{std::uint16_t}、x86 word \\
\code{dword ptr} & 4 & \code{int}、\code{float}、\code{std::uint32_t} \\
\code{qword ptr} & 8 & \code{long}、指针、\code{double}、\code{std::uint64_t} \\
\code{xmmword ptr} & 16 & 128-bit SIMD memory operand \\
\code{ymmword ptr} & 32 & 256-bit AVX memory operand \\
\bottomrule
\end{tabular}
\end{center}

例如：

\begin{lstlisting}
movzx eax, byte ptr [rdi]
mov   eax, dword ptr [rdi]
mov   rax, qword ptr [rdi]
movss xmm0, dword ptr [rdi]
movsd xmm0, qword ptr [rdi]
\end{lstlisting}

这些指令即使地址相同，访问宽度和解释规则也不同：

\begin{itemize}
  \item \instruction{movzx eax, byte ptr [rdi]} 读取 1 字节并零扩展。
  \item \instruction{mov eax, dword ptr [rdi]} 读取 4 字节整数到 \register{eax}。
  \item \instruction{mov rax, qword ptr [rdi]} 读取 8 字节整数或指针。
  \item \instruction{movss xmm0, dword ptr [rdi]} 读取单精度浮点标量。
  \item \instruction{movsd xmm0, qword ptr [rdi]} 读取双精度浮点标量。
\end{itemize}

\begin{keyidea}
内存地址只说明从哪里开始，操作数宽度说明读写多少字节，指令语义和寄存器类型说明如何解释这些字节。三者缺一不可。
\end{keyidea}

C++ 类型信息在汇编中大多已经消失。你看到 \code{dword ptr [rdi]}，只能知道访问 4 字节，不能仅凭这一点判断它一定是 \code{int}。它也可能是 \code{float}，可能是无符号整数，可能是结构体中四字节字段，甚至可能只是某种位模式。真正区分整数和浮点的，往往是使用的寄存器类别和指令语义：通用寄存器上的 \instruction{mov}、\instruction{add} 更像整数或地址语义；XMM 寄存器上的 \instruction{movss}、\instruction{addss} 更像单精度浮点语义。

访问宽度也不等于对象大小。有时编译器只读取对象的一部分，例如读取一个字节字段或布尔值；有时会使用更宽的加载再用掩码取出需要的位；有时为了向量化会一次读取多个元素；有时结构体拷贝会变成若干宽 load/store。后续分析对象布局时，必须把“这条指令访问了多少字节”和“源语言对象有多大”分开。

宽度还直接影响正确性。读取 1 字节后如果要变成 32 位整数，就必须决定高位如何处理：零扩展还是符号扩展。读取 4 字节到 \register{eax} 会清零 \register{rax} 高位，但读取到 \register{al} 不会。很多 C/C++ 类型转换、数组元素读取和字符处理的汇编差异，都来自这个问题。下一章讲整数指令时会系统展开 \instruction{movzx}、\instruction{movsx} 和各种扩展规则，本章先把“宽度必须标清”作为基本习惯。

\section{浮点和 SIMD 基础：看到 xmm0 时先问什么}

第一册不会把 SIMD 优化讲成第二册那样深，但读 x86-64 汇编时必须先认识浮点和向量指令。原因很简单：现代编译器即使编译很普通的 C++，也可能用 SSE/AVX 指令处理浮点、拷贝、小结构体、自动向量化循环或库函数调用。若读者只认识 \register{rax}、\register{rdi} 和整数 \instruction{add}，看到 \register{xmm0} 或 \instruction{vaddps} 就会把真实热路径跳过去。

\subsection{XMM、YMM、ZMM：同一组向量寄存器的不同宽度视图}

x86-64 上常见的浮点和 SIMD 寄存器包括：

\begin{center}
\small
\begin{tabular}{p{0.16\linewidth}p{0.18\linewidth}p{0.46\linewidth}}
\toprule
寄存器名 & 常见宽度 & 读汇编时的重点 \\
\midrule
\register{xmm0}--\register{xmm15} & 128 bit & SSE 和标量浮点常用；System V ABI 中浮点参数和返回值也常走 \register{xmm0} 起 \\
\register{ymm0}--\register{ymm15} & 256 bit & AVX/AVX2 常用；低 128 bit 与对应 XMM 视图重叠 \\
\register{zmm0}--\register{zmm31} & 512 bit & AVX-512 常用；还会配合 \register{k0}--\register{k7} mask 寄存器 \\
\bottomrule
\end{tabular}
\end{center}

这些名字不是三套互不相干的寄存器。可以先把它们理解成同一条物理/架构向量寄存器的不同宽度视图：\register{xmm0} 是低 128 位，\register{ymm0} 包含低 256 位，\register{zmm0} 包含低 512 位。实际 CPU 内部实现会更复杂，但汇编阅读的第一层先抓住这个视图关系。

\begin{warningbox}
不要把 \register{xmm0} 简单理解成“一个 float”。它可以装一个标量 \code{float}，也可以装两个 \code{double}、四个 \code{float}、十六个字节，甚至只是被当作位模式搬运。解释来自指令后缀和上下文。
\end{warningbox}

\subsection{ss、sd、ps、pd：后缀就是 lane 语义}

浮点指令名里的后缀通常直接告诉你元素类型和是否 packed：

\begin{center}
\small
\begin{tabular}{p{0.15\linewidth}p{0.26\linewidth}p{0.43\linewidth}}
\toprule
后缀 & 含义 & 例子 \\
\midrule
\code{ss} & scalar single，一个低位 \code{float} & \instruction{addss xmm0, xmm1} 只对低 32 位单精度元素做加法 \\
\code{sd} & scalar double，一个低位 \code{double} & \instruction{addsd xmm0, xmm1} 只对低 64 位双精度元素做加法 \\
\code{ps} & packed single，多个 \code{float} lane & \instruction{addps xmm0, xmm1} 对 4 个 \code{float} lane 分别加法 \\
\code{pd} & packed double，多个 \code{double} lane & \instruction{addpd xmm0, xmm1} 对 2 个 \code{double} lane 分别加法 \\
\bottomrule
\end{tabular}
\end{center}

同样是加法，下面两条含义完全不同：

\begin{lstlisting}
addss xmm0, xmm1    ; one scalar f32 in the low lane
addps xmm0, xmm1    ; four packed f32 lanes in xmm0/xmm1
\end{lstlisting}

\instruction{addss} 不是“慢一点的 \instruction{addps}”，它的语义就是只更新低位单精度标量。高位部分如何保留或清零，还要看是 legacy SSE 编码还是 VEX/EVEX 编码，以及具体指令规则。第一层读法不要猜：看到标量后缀，就在笔记里写“低 lane 是业务值，高 lane 不一定有业务意义”。

\subsection{SSE 二操作数和 AVX 三操作数：目的寄存器是否被破坏}

legacy SSE 指令通常是二操作数形式，目的操作数同时也是输入：

\begin{lstlisting}
addps xmm0, xmm1    ; xmm0 = xmm0 + xmm1
\end{lstlisting}

AVX 的 VEX 编码通常提供三操作数形式：

\begin{lstlisting}
vaddps ymm0, ymm1, ymm2    ; ymm0 = ymm1 + ymm2
\end{lstlisting}

这个差异对读数据流很重要。第一条会消耗旧 \register{xmm0}，再把结果写回 \register{xmm0}；第二条把 \register{ymm1} 和 \register{ymm2} 作为输入，把结果写到 \register{ymm0}。如果你把二者都注释成“向量相加”，就漏掉了依赖链。

\begin{keyidea}
读 SIMD 指令时，先写 lane 类型，再写 lane 数，再写目的寄存器是否也是输入。性能分析之前，先把这三个语义事实写清楚。
\end{keyidea}

\subsection{最小 SIMD 词典：load、算术、shuffle、reduce、mask}

第一册只要求建立词典，不要求优化。看到 SIMD 反汇编时，先把指令粗分成几类：

\begin{center}
\small
\begin{tabular}{p{0.18\linewidth}p{0.32\linewidth}p{0.34\linewidth}}
\toprule
类别 & 常见形态 & 阅读重点 \\
\midrule
load/store & \instruction{movups}、\instruction{movaps}、\instruction{vmovdqu}、\instruction{vmovaps} & 地址、宽度、是否要求对齐、是否跨 cache line 或页 \\
lane 算术 & \instruction{addps}、\instruction{mulps}、\instruction{vaddps}、\instruction{vfmadd231ps} & 每个 lane 独立计算，FMA 通常按乘加两个浮点操作计 \\
shuffle/permute & \instruction{shufps}、\instruction{unpcklps}、\instruction{vpermps} & 改变 lane 位置，常比普通加法更难吞吐 \\
horizontal reduce & \instruction{haddps} 或若干 shuffle + add & 把多个 lane 合成一个标量，常是 dot/sum 的尾部成本 \\
mask/tail & AVX-512 \register{k} mask，或 AVX2 blend/标量尾部 & 处理长度不是向量宽度整数倍的尾部元素 \\
\bottomrule
\end{tabular}
\end{center}

例如 \instruction{vfmadd231ps ymm0, ymm1, ymm2} 通常表示 packed single FMA，但 \code{132}、\code{213}、\code{231} 的操作数顺序不适合凭名字猜。报告里应直接引用反汇编器注释或查手册，并写成明确公式。若工具显示：

\begin{lstlisting}
vfmadd231ps ymm0, ymm1, ymm2   # ymm0 = (ymm1 * ymm2) + ymm0
\end{lstlisting}

就把这条公式写进分析。不要只写“FMA”，因为不同 FMA 变体的累加寄存器位置不同，错误理解会直接把 dot product 数据流看错。

\subsection{对齐、上半寄存器和浮点语义边界}

SIMD 入门有三个常见坑。

第一，对齐不是装饰。某些 aligned load/store 形式要求地址对齐，例如 \instruction{movaps} 或 \instruction{vmovaps} 的内存操作数；不满足时可能 fault，或者至少说明接口合同不清。对应的 unaligned 形式如 \instruction{movups}、\instruction{vmovups} 更宽容，但性能仍可能受跨 cache line、跨页和地址生成影响。报告不能只写“加载 xmmword”，还要写具体指令和地址对齐前提。

第二，legacy SSE 和 AVX 混用可能涉及 upper-lane 状态。简化地说，旧 SSE 指令经常只写 XMM 低 128 位视图，而 AVX/VEX 指令有更明确的向量宽度和上半部分处理规则。真实性能问题中还会遇到 SSE/AVX transition penalty、\instruction{vzeroupper} 等内容；第一册只要求建立边界：看到 \register{xmm} 和 \register{ymm} 混用时，不要只看低 128 位结果，还要记录这是 legacy SSE 还是 VEX 编码。

第三，浮点不是实数。NaN、Inf、signed zero、舍入模式、非规格化数、MXCSR 中的 FTZ/DAZ、\code{-ffast-math} 都可能改变结果边界。第一册的汇编阅读只要求你知道：\instruction{addss}、\instruction{mulps} 和 \instruction{vfmadd} 不是数学实数运算，它们遵守浮点环境和指令语义。做 benchmark 或优化时，必须有误差模型。

\subsection{浮点/SIMD 反汇编报告最少字段}

读到浮点或 SIMD 指令时，报告至少写：

\begin{itemize}
  \item 寄存器宽度：XMM、YMM 还是 ZMM。
  \item lane 类型和 lane 数：例如 \code{ymm} 上的 \code{ps} 是 8 个 \code{float} lane。
  \item 指令形态：legacy SSE 二操作数，还是 VEX/EVEX 三操作数。
  \item load/store 宽度和地址公式。
  \item 对齐要求和尾部处理方式。
  \item 浮点误差边界：bit-exact、容差、是否允许 FMA 或 fast-math。
\end{itemize}

例如：

\begin{lstlisting}[numbers=none,caption={SIMD 指令阅读记录示例}]
instruction:
  vaddps ymm0, ymm1, ymm2
facts:
  vector width = 256 bit
  lane type = f32
  lane count = 8
  data flow = ymm0 receives ymm1 + ymm2
boundary:
  This proves packed f32 lane-wise add.
  It does not prove the source loop had no aliasing problem.
\end{lstlisting}

\subsection{IEEE 754：浮点数是有格式的位串，不是实数}

读到 \instruction{addss}、\instruction{addsd}、\instruction{mulps}、\instruction{vfmadd231ps} 时，不能把它们当成数学实数运算。x86-64 上常见的 \code{float} 和 \code{double} 通常使用 IEEE 754 二进制浮点格式。它们首先是固定宽度位串，然后才按符号位、指数和尾数解释成数值。

\begin{center}
\small
\begin{tabular}{p{0.16\linewidth}p{0.16\linewidth}p{0.16\linewidth}p{0.16\linewidth}p{0.22\linewidth}}
\toprule
类型 & 总位数 & 符号位 & 指数字段 & 尾数字段 \\
\midrule
\code{float} & 32 & 1 & 8 & 23，隐藏 leading 1 只对规格化数成立 \\
\code{double} & 64 & 1 & 11 & 52，隐藏 leading 1 只对规格化数成立 \\
\bottomrule
\end{tabular}
\end{center}

这个格式直接解释了几个工程现象。第一，浮点有正零和负零。它们比较时通常相等，但某些运算、打印、倒数和符号传播会区分。第二，浮点有无穷大，除以零、溢出或某些库函数可能产生 \code{+inf} 或 \code{-inf}。第三，浮点有 NaN，表示不是普通数值的结果，例如 \code{0.0 / 0.0} 或无效运算。第四，浮点还有非规格化数，也叫 subnormal 或 denormal，用来表示非常接近 0 的小数，但可能触发慢路径或被运行环境刷成 0。

\begin{warningbox}
不要把 \code{float} 解释成“32 位实数”。它是 32 位 IEEE 754 编码。有限精度、特殊值、舍入模式和浮点环境都是语义的一部分。
\end{warningbox}

\subsection{舍入模式：为什么浮点加法不满足普通代数直觉}

IEEE 754 运算的结果通常要舍入到目标格式。默认舍入模式常见为 round to nearest, ties to even，也就是最接近可表示值，正好居中时选尾数最低位为偶数的那个结果。除此之外，浮点环境还可能选择向零、向正无穷、向负无穷舍入。

舍入会破坏很多数学直觉：

\begin{lstlisting}[language=C++]
float a = 1.0e20f;
float b = -1.0e20f;
float c = 3.14f;

float x = (a + b) + c;
float y = a + (b + c);
\end{lstlisting}

数学实数里两者相等；浮点里未必。因为 \code{b + c} 可能先把 \code{c} 舍掉，后续再和 \code{a} 相加就得到不同结果。优化器、向量化和 FMA 都会遇到这个边界：如果编译选项允许重关联或 fast-math，编译器可能把表达式改成数学上等价但 IEEE 结果不同的形态；如果要求严格浮点语义，编译器必须更保守。

FMA 是另一个典型例子。表达式：

\begin{lstlisting}[language=C++]
float r = a * b + c;
\end{lstlisting}

可能生成分离的乘法和加法，也可能生成 fused multiply-add。分离版本先把 \code{a * b} 舍入成 \code{float}，再加 \code{c} 并再次舍入；FMA 通常在更高精度中完成乘加，只舍入一次。因此 FMA 往往更准，也可能和非 FMA 结果最后几位不同。报告中若看到 \instruction{vfmadd231ps}，不能只写“等价于乘法再加法”，而要写清它是一次 fused operation。

\subsection{MXCSR：SSE/AVX 浮点环境的关键状态}

整数汇编常常只需要看通用寄存器和 flags；浮点/SIMD 汇编还要知道 MXCSR。MXCSR 是 SSE/AVX 浮点环境中的控制和状态寄存器，包含舍入控制、异常标志、异常 mask，以及常见实现中的 FTZ/DAZ 控制。

\begin{center}
\small
\begin{tabular}{p{0.22\linewidth}p{0.34\linewidth}p{0.28\linewidth}}
\toprule
字段或概念 & 含义 & 工程影响 \\
\midrule
rounding control & 控制浮点结果如何舍入 & 可改变边界结果和测试 oracle \\
exception flags & 记录 invalid、divide-by-zero、overflow、underflow、inexact 等事件 & 可用于诊断数值问题，但常被忽略 \\
exception masks & 控制异常是置标志还是触发 trap & 默认多数环境会 mask，程序继续运行 \\
FTZ & flush-to-zero，把某些 underflow 结果刷成 0 & 可避免 subnormal 慢路径，但改变严格结果 \\
DAZ & denormals-are-zero，把输入 subnormal 当 0 & 可提升速度，但改变输入语义 \\
\bottomrule
\end{tabular}
\end{center}

这解释了为什么“同一段浮点代码”在不同运行环境下可能不完全一致。一个 HPC 或 AI kernel 可能为了速度打开 FTZ/DAZ；一个数值库可能要求更严格的 IEEE 行为；一个测试如果没有记录 MXCSR，就可能在某台机器通过，在另一台机器最后几位不同或 denormal 路径耗时异常。

第一册不要求手写 MXCSR 管理代码，但报告要知道怎么提问：

\begin{lstlisting}[numbers=none,caption={浮点环境报告字段}]
floating_environment:
  target_format: f32 / f64
  rounding_mode:
  fma_allowed:
  fast_math_flags:
  mxcsr:
    ftz:
    daz:
    exception_masks:
    exception_flags_before_after:
\end{lstlisting}

\subsection{NaN、无穷大和 unordered 比较}

NaN 是浮点错误分析里最容易被低估的值。NaN 与任何值比较都不满足普通全序关系，包括它自己。常见现象是：

\begin{lstlisting}[language=C++]
double x = std::numeric_limits<double>::quiet_NaN();
bool a = (x == x);   // usually false
bool b = (x <  1.0); // false
bool c = (x >  1.0); // false
\end{lstlisting}

x86 SSE 比较指令会把 NaN 情况归为 unordered。你可能看到 \instruction{ucomiss}、\instruction{ucomisd}、\instruction{comiss}、\instruction{comisd} 后接条件跳转。读这类代码时不能只套整数比较直觉，因为 NaN 会影响 flags 组合和跳转选择。编译器为了实现 \code{std::isnan}、\code{x == x}、\code{x < y} 等语义，可能生成看起来不直观的条件码序列。

无穷大也不是普通大数。有限数乘到溢出可能变成 \code{inf}；\code{1.0 / +0.0} 可能得到 \code{+inf}；\code{1.0 / -0.0} 可能得到 \code{-inf}。如果优化把某个分母认为“不可能为 0”，那通常需要语言前提、范围分析或 fast-math 许可。没有这些前提，直接把除法、比较或异常路径删除可能改变 IEEE 语义。

\subsection{denormal/subnormal：正确性值和性能慢路径}

非规格化数用于表示比最小规格化数更接近 0 的值。它们让浮点下溢附近不会突然断崖式变成 0，这叫 gradual underflow。但硬件处理 subnormal 可能明显慢，尤其是在旧 CPU 或某些指令路径上。许多实时音频、HPC 和推理 kernel 会选择打开 FTZ/DAZ，避免小值进入 denormal 慢路径。

这带来一个工程取舍：

\begin{itemize}
  \item 严格 IEEE 路径保留 subnormal，结果更符合标准语义。
  \item FTZ/DAZ 路径把极小值当 0，通常更快，但结果边界不同。
  \item benchmark 如果输入包含大量极小值，可能测到的不是普通算术吞吐，而是 denormal 处理成本。
\end{itemize}

AI Infra 中常见的 fp32/fp16/bf16 kernel 也有类似问题。很多模型计算不要求 bit-exact，但必须写清误差边界和浮点环境。否则一个 kernel 在随机常规输入上很快，在接近 0 的 adversarial 输入上突然变慢或结果不一致，报告会无法解释。

\subsection{搬运位模式和数值转换不是一回事}

XMM/YMM 指令里有很多看起来都像“移动”的操作，但语义不同。下面这些要分清：

\begin{lstlisting}
movss  xmm0, dword ptr [rdi]   ; load low f32 lane as bits
movsd  xmm0, qword ptr [rdi]   ; load low f64 lane as bits
movd   xmm0, eax               ; move 32 raw bits between GPR and XMM
cvtss2sd xmm0, xmm1            ; convert f32 value to f64 value
cvtsi2sd xmm0, edi             ; convert integer value to f64 value
\end{lstlisting}

\instruction{movd} 搬的是位模式，不把整数数值“变成浮点数”。\instruction{cvtsi2sd} 才是整数到 double 的数值转换。反过来，\instruction{cvttsd2si} 这类指令把浮点转换成整数，并涉及舍入或截断、溢出和异常边界。读汇编时如果把 bit move 和 numeric conversion 混掉，就会把 ABI glue、类型转换和优化序列全部看错。

\subsection{XMM/YMM 高位：标量低 lane 不代表高 lane 有意义}

标量浮点指令通常只把业务值放在低 lane。高 lane 可能保留旧值，也可能被清零，具体取决于 legacy SSE、VEX/EVEX 编码和指令规则。比如 \instruction{addss} 的核心业务值是低 32 位标量；\instruction{vaddss} 的三操作数形态还涉及高位从哪个源寄存器继承。阅读时应避免把整个 \register{xmm0} 都解释成一个“干净对象”。

常见清零 idiom：

\begin{lstlisting}
vxorps xmm0, xmm0, xmm0
xorps  xmm0, xmm0
\end{lstlisting}

它们通常用于把寄存器置零，但在 AVX/SSE 混合代码中，VEX 编码的零 idiom 还可能帮助处理上半寄存器状态。函数从 AVX 代码返回到可能使用 legacy SSE 的调用者前，编译器或库函数常见 \instruction{vzeroupper}，用于避免某些微架构上的 SSE/AVX 过渡惩罚。第一册不要求深入性能细节，但报告要能识别这不是“无意义清零”。

\subsection{浮点反例：为什么 bit-exact 测试经常不合适}

浮点优化的正确性不能简单套整数 bit-exact 习惯。下面几个变化都可能让最后几位不同：

\begin{itemize}
  \item 标量循环改成 SIMD reduction，改变求和顺序。
  \item 分离乘加改成 FMA，改变舍入次数。
  \item 开启 \code{-ffast-math}，允许重关联、忽略部分 NaN/Inf 语义。
  \item 打开 FTZ/DAZ，改变 subnormal 输入或输出。
  \item 使用不同 libm 实现或目标 ISA，改变近似函数结果。
\end{itemize}

因此，浮点测试报告要写清：本函数要求 bit-exact，还是允许绝对误差/相对误差/ULP 误差；是否允许 FMA；是否允许 fast-math；是否测试 NaN、Inf、signed zero、subnormal 和极大值。没有这些字段，只说“结果差一点”或“结果一致”都不够专业。

\section{内存寻址模式：地址公式如何写进指令}

x86-64 最常见的内存操作数形式是：

\begin{lstlisting}
[base + index * scale + displacement]
\end{lstlisting}

其中 \code{scale} 只能是 1、2、4、8。不是所有部分都必须出现：

\begin{center}
\begin{tabular}{p{0.35\linewidth}p{0.48\linewidth}}
\toprule
形式 & 含义 \\
\midrule
\code{[rdi]} & 地址是 \register{rdi} \\
\code{[rdi + 8]} & 地址是 \register{rdi} 加 8 字节 \\
\code{[rdi + rsi*4]} & 地址是 \register{rdi} 加 \register{rsi} 乘 4 \\
\code{[rdi + rsi*4 + 12]} & 基址、索引缩放、常量偏移同时出现 \\
\code{[rip + symbol]} & RIP-relative，常用于全局对象或常量池 \\
\bottomrule
\end{tabular}
\end{center}

\subsection{案例：数组 load}

C++：

\begin{lstlisting}[language=C++]
extern "C" int load_i32(int const* p, int i)
{
    return p[i];
}
\end{lstlisting}

可能汇编：

\begin{lstlisting}
movsxd rsi, esi
mov    eax, dword ptr [rdi + rsi*4]
ret
\end{lstlisting}

假设执行前：

\begin{center}
\begin{tabular}{lll}
\toprule
寄存器 & 值 & 含义 \\
\midrule
\register{rdi} & \code{0x1000} & \code{p[0]} 地址 \\
\register{esi} & \code{3} & 下标 \code{i} \\
\bottomrule
\end{tabular}
\end{center}

执行：

\begin{lstlisting}
movsxd rsi, esi
    rsi = sign_extend_64(3)
        = 3

mov eax, dword ptr [rdi + rsi*4]
    effective_address = 0x1000 + 3 * 4
                      = 0x100c
    load 4 bytes from 0x100c..0x100f into eax
\end{lstlisting}

\subsection{案例：结构体字段}

C++：

\begin{lstlisting}[language=C++]
struct S {
    int a;
    double b;
    int c;
};

extern "C" int get_c(S const* p)
{
    return p->c;
}
\end{lstlisting}

常见布局：

\begin{center}
\begin{tabular}{llll}
\toprule
字节范围 & 内容 & 大小 & 说明 \\
\midrule
0..3 & \code{a} & 4 & 字段 \\
4..7 & padding & 4 & 让 \code{b} 8 字节对齐 \\
8..15 & \code{b} & 8 & 字段 \\
16..19 & \code{c} & 4 & 字段 \\
20..23 & tail padding & 4 & 让 \code{sizeof(S)} 为 8 的倍数 \\
\bottomrule
\end{tabular}
\end{center}

可能汇编：

\begin{lstlisting}
mov eax, dword ptr [rdi + 16]
ret
\end{lstlisting}

这里 \code{+ 16} 不是魔法数字，而是 \code{offset(S, c)}：

\begin{lstlisting}
addr(p->c) = addr(*p) + offset(S, c)
           = rdi      + 16
\end{lstlisting}

\subsection{案例：结构体数组字段}

C++：

\begin{lstlisting}[language=C++]
struct Pair {
    int a;
    int b;
};

extern "C" int get_b(Pair const* p, int i)
{
    return p[i].b;
}
\end{lstlisting}

公式：

\begin{lstlisting}
sizeof(Pair)    = 8
offset(Pair, b) = 4
addr(p[i].b)    = base + i * 8 + 4
\end{lstlisting}

可能汇编：

\begin{lstlisting}
movsxd rsi, esi
mov    eax, dword ptr [rdi + rsi*8 + 4]
ret
\end{lstlisting}

如果 \register{rdi} 是 \code{0x2000}，\register{rsi} 是 \code{5}：

\begin{lstlisting}
effective_address = 0x2000 + 5 * 8 + 4
                  = 0x202c
load 4 bytes from 0x202c..0x202f
\end{lstlisting}

\subsection{地址公式要同时解释三件事}

一个内存操作数的地址公式至少有三层含义。第一层是机器层的有效地址，也就是寄存器、缩放和位移如何相加。这一层必须严格按照汇编文本写出来，不能省略。第二层是布局层的来源，也就是为什么缩放是 4、8 或其他形式，为什么常量偏移是 8、12、16。它可能来自元素大小、字段偏移、栈槽位置或表项大小。第三层才是源码层解释，也就是它可能对应 \code{p[i]}、\code{obj->field}、\code{items[i].count} 或某个临时 spill。

这三层的证据强度不同。机器层是确定的；布局层通常可以通过类型大小、对齐和相邻访问推导；源码层则需要结合函数签名和编译器优化。高质量汇编报告应该把三层分开写。例如不要只说“访问数组第 i 个元素”，而要写“有效地址为 \register{rdi} 加 \register{rsi} 乘 4；若 \register{rdi} 是 \code{int} 数组基址且 \register{rsi} 是下标，则它对应 \code{p[i]}；访问宽度为 4 字节”。

地址公式还能暴露数据布局问题。结构体数组字段访问常常出现“索引乘结构体大小再加字段偏移”；二维数组访问会出现“行号乘行跨度再加列偏移”；栈上局部变量常常表现为 \register{rsp} 或 \register{rbp} 加减常量。以后读更复杂代码时，你会靠这些公式恢复编译器看到的数据布局。第 6 章的目标不是背所有模式，而是训练你看到地址公式就能问出正确问题。

\begin{warningbox}
不要把 displacement 叫作“随机常数”。汇编里的常量偏移通常有来源：字段偏移、数组元素偏移、栈槽位置、全局对象地址差、跳转目标位移或对齐填充。找不到来源时，应说“暂未解释”，不要编造一个源码形状。
\end{warningbox}

\section{RIP-relative：全局对象和常量池的常见形态}

x86-64 代码里经常出现 \code{[rip + ...]} 形式：

\begin{lstlisting}
mov eax, dword ptr [rip + global_value]
lea rdi, [rip + message]
\end{lstlisting}

这里的 \register{rip} 不是普通基址寄存器用法，而是 RIP-relative addressing。它用“下一条指令地址附近的相对位移”来定位全局对象、字符串字面量、常量池或跳转表。这种形式对位置无关代码非常重要，因为代码被加载到不同基址时，相对距离可以由链接器和加载器处理。

考虑：

\begin{lstlisting}[language=C++]
extern int global_value;

extern "C" int read_global()
{
    return global_value;
}
\end{lstlisting}

目标文件阶段可能看到类似：

\begin{lstlisting}
mov eax, dword ptr [rip + 0x0]
        R_X86_64_PC32 global_value-0x4
ret
\end{lstlisting}

这不是说全局变量地址就是 0。目标文件还不知道 \code{global_value} 最终放在哪里，所以在指令里留下占位位移，并在 relocation 中记录“链接时把这里修成相对 \code{global_value} 的位移”。链接后再看最终可执行文件，位移可能变成具体数值。

RIP-relative 常见于：

\begin{itemize}
  \item 读取只读常量或字符串地址。
  \item 访问全局变量。
  \item 访问 GOT 表项。
  \item 访问跳转表或常量表。
  \item 生成位置无关代码中的相对地址。
\end{itemize}

读这种地址时要小心两点。第一，\code{[rip + disp]} 的 \register{rip} 通常指向下一条指令，而不是当前指令起始地址。第二，反汇编器可能把目标符号名显示出来，也可能只显示数值位移；是否显示符号取决于符号表、重定位信息和工具选项。

\begin{mentalmodel}
普通 \code{[rdi + rsi*4]} 常常来自对象或数组地址；\code{[rip + disp]} 常常来自“代码附近的全局/只读/链接时符号”。看到 RIP-relative，优先想到链接、位置无关代码和全局数据。
\end{mentalmodel}

\section{重定位：目标文件里尚未完成的地址}

汇编阅读不能只看最终指令文本，还要知道哪些地址在目标文件阶段尚未确定。重定位就是工具链记录“这里以后需要修补”的机制。

例如一个文件调用另一个文件里的函数：

\begin{lstlisting}[language=C++]
extern "C" int callee(int);

extern "C" int caller(int x)
{
    return callee(x + 1);
}
\end{lstlisting}

编译成目标文件后，\code{caller.o} 中可能有一条 \instruction{call}，但它还不知道 \code{callee} 的最终地址。反汇编可能显示：

\begin{lstlisting}
83 c7 01              add edi, 1
e8 00 00 00 00        call 0x8
        R_X86_64_PLT32 callee-0x4
c3                    ret
\end{lstlisting}

\code{e8 00 00 00 00} 中的 4 字节位移暂时是 0；重定位项告诉链接器，这里要修成到 \code{callee} 的相对位移，可能经过 PLT。链接后，这条 call 的位移才变成最终值。

因此，目标文件反汇编里的地址有三类：

\begin{itemize}
  \item 已经确定的本地偏移，例如同一函数内跳转到局部标签。
  \item 需要链接器修补的外部符号引用。
  \item 需要动态链接器或加载器参与的动态重定位。
\end{itemize}

读报告时，如果作者只贴目标文件中的 \instruction{call 0x0} 或 \instruction{mov [rip+0x0]}，却不贴 relocation，就不完整。因为这些 0 很可能只是占位，不是实际目标地址。正确证据应同时包含 \code{objdump -drwC -Mintel} 的指令和重定位行，或者直接分析链接后的可执行文件。

\begin{keyidea}
目标文件不是最终程序。看到外部函数、全局对象、RIP-relative 占位或 \instruction{call} 目标异常时，先查重定位，再下结论。
\end{keyidea}

\section{如何组合 \code{objdump}、\code{readelf} 和 \code{nm}}

单个工具很少给出完整答案。读目标文件时，\code{objdump} 擅长显示指令、机器码字节和重定位旁注。\code{readelf} 擅长显示 ELF section、符号表和重定位表的结构化信息。\code{nm} 擅长快速列出符号定义和未定义引用。高质量分析应根据问题组合工具，而不是把某个工具输出当作全部真相。

\begin{center}
\begin{tabular}{p{0.22\linewidth}p{0.30\linewidth}p{0.34\linewidth}}
\toprule
工具 & 常用命令 & 回答的问题 \\
\midrule
\code{objdump} & \code{objdump -drwC ...} & 指令、字节、反汇编、重定位旁注 \\
\code{readelf} & \code{readelf -S -s -r file.o} & section、符号表、重定位表 \\
\code{nm} & \code{nm -C file.o} & 哪些符号已定义，哪些仍未定义 \\
\code{llvm-objdump} & \code{llvm-objdump -dr ...} & LLVM 工具链下的反汇编和重定位 \\
\bottomrule
\end{tabular}
\end{center}

完整命令可以写在代码块中，避免表格因为长选项变得难读：

\begin{lstlisting}[language=bash]
objdump -drwC -Mintel file.o
llvm-objdump -dr --x86-asm-syntax=intel file.o
\end{lstlisting}

比如你看到目标文件中有：

\begin{lstlisting}
e8 00 00 00 00        call 0x5
        R_X86_64_PLT32 printf-0x4
\end{lstlisting}

这时不能说“它调用地址 \code{0x5}”。\code{0x5} 是当前目标文件里占位位移导致的显示形式；重定位行才说明这里引用 \code{printf}。如果再运行 \code{nm -C file.o}，可能看到：

\begin{lstlisting}
                 U printf
\end{lstlisting}

这说明 \code{printf} 在当前目标文件中未定义，链接器需要从 C 运行时库或动态库中解析它。若运行 \code{readelf -r file.o}，会看到重定位表中有对应条目，说明哪个 offset 需要修补、重定位类型是什么、关联哪个符号。

如果问题是“函数里实际有哪些指令”，优先看 \code{objdump}。如果问题是“为什么 call 目标是 0”，必须看重定位。问题是“这个外部名字是否已经定义”，看 \code{nm}。问题是“这个常量放在哪个 section，section 是否可写”，看 \code{readelf -S}。问题是“符号大小、绑定和类型是什么”，看 \code{readelf -s} 或 \code{objdump -t}。同一个二进制事实经常要从多个角度交叉确认。

\subsection{目标文件、可执行文件和运行时地址}

同一个函数在三个阶段会有不同“地址语境”。在目标文件中，地址通常是 section 内偏移，很多外部地址还没解析；在链接后的可执行文件或共享库中，链接器已经给 section 分配虚拟地址，许多符号引用已经确定或变成动态链接入口；在运行中的进程里，PIE、ASLR 和动态加载会让实际装载基址变化。

因此下面三句话不是一回事：

\begin{itemize}
  \item “目标文件中 \code{add2} 位于 \code{.text} offset 0。”
  \item “链接后可执行文件中 \code{add2} 的虚拟地址是 \code{0x401140}。”
  \item “当前进程这次运行中 \code{add2} 实际映射在某个运行时地址。”
\end{itemize}

调试器、perf、反汇编和链接器 map 文件可能使用不同地址语境。若报告中把它们混在一起，就会出现“objdump 里是 0，但 GDB 里不是 0”“上次运行地址和这次不一样”这类困惑。正确做法是每次写地址都标明阶段：目标文件 offset、链接后虚拟地址，还是运行时地址。

\subsection{不要忽略工具选项}

工具选项会改变你看到的信息。只用 \code{objdump -d}，可能看不到重定位；加上 \code{-r} 才显示 relocation；加上 \code{-w} 可以减少长行截断；加上 \code{-C} 可以反修饰 C++ 符号名；加上 \code{-Mintel} 可以显示 Intel 语法。不同平台的 \code{objdump} 和 \code{llvm-objdump} 选项略有差异，报告里应写明使用的命令。

这不是形式主义。一个只贴 \code{objdump -d} 的报告可能漏掉所有外部符号修补信息；一个没有 \code{-C} 的报告可能把 C++ 符号名写成难读的 mangled name；一个没有说明 Intel/AT\&T 语法的报告可能把操作数方向解释反。工具命令就是证据链的一部分。

\begin{keyidea}
反汇编不是“把二进制还原成源码”。它只是把字节解释成指令文本，并尽量附上符号和重定位信息。真正的分析要同时看指令、section、符号、重定位、编译命令和 ABI。
\end{keyidea}

\section{外部函数、PLT 和 GOT 的第一册认识}

当代码调用外部函数或访问动态库中的全局对象时，反汇编中可能出现 \code{@PLT}、\code{@GOTPCREL}、GOT load、PLT stub 等字样。完整动态链接机制比较大，后续可以单独成章；本章先建立阅读边界，避免把这些名字当作普通业务函数或普通数组访问。

例如：

\begin{lstlisting}
call printf@PLT
mov rax, qword ptr [rip + global@GOTPCREL]
\end{lstlisting}

\code{PLT} 通常与过程链接表有关，用来支持对动态库函数的间接解析和跳转；\code{GOT} 通常与全局偏移表有关，用来保存位置无关代码访问外部符号所需的地址。你不需要在第 6 章掌握动态链接全部细节，但看到这些标记时，至少要知道：这里不是一个普通的局部 \instruction{call} 或普通的数组 load，它经过链接器、动态链接器和位置无关代码约定。

外部函数调用在不同构建方式下形态会变。静态链接、动态链接、PIE、非 PIE、默认可见性、隐藏可见性、链接时优化、是否允许符号插桩或 interposition，都可能影响最终指令。有时你看到 \instruction{call printf@PLT}；有时看到对 GOT 表项的间接跳转；有时链接器能把调用放松成更直接的相对 call。报告里不要把某一种形态当作所有 Linux 程序的唯一形态。

从学习路径看，第一册只要求你做到三点。第一，看到 \code{@PLT} 或 \code{@GOTPCREL} 时意识到这是链接和动态加载相关机制。第二，能通过 \code{objdump -drwC} 和 \code{readelf -r} 找到对应重定位。第三，不把 PLT/GOT 开销直接归因成业务算法瓶颈，除非有运行时采样或计数证据证明它确实在热路径中频繁发生。

\begin{warningbox}
PLT/GOT 是二进制链接边界的一部分，不是 C++ 源码里的普通函数体。读外部调用时，要把“业务函数语义”和“动态链接入口机制”分开。
\end{warningbox}

\section{栈上的地址公式：\register{rsp}、\register{rbp} 和栈槽}

汇编里除了数组和结构体地址，还经常看到以 \register{rsp} 或 \register{rbp} 为基址的内存操作数：

\begin{lstlisting}
mov dword ptr [rsp + 12], edi
mov eax, dword ptr [rbp - 4]
mov qword ptr [rsp + 24], rax
\end{lstlisting}

这些通常和栈帧、局部对象、保存寄存器、溢出临时值或调用参数有关。初学者容易把每个栈槽都对应成一个源码局部变量，这在 \code{-O0} 下有时近似成立，在优化构建中经常不成立。

栈地址阅读要先回答两个问题。

第一，当前函数是否使用 frame pointer。如果函数序言里有：

\begin{lstlisting}
push rbp
mov rbp, rsp
\end{lstlisting}

那么 \register{rbp} 常被用作稳定基准，局部变量可能位于 \code{[rbp - offset]}，调用者相关数据或栈上传参可能位于 \code{[rbp + offset]}。如果省略 frame pointer，编译器通常用 \register{rsp} 加偏移访问栈槽，但 \register{rsp} 可能随 push、call、栈空间调整而变化，阅读时要追踪它的当前值。

第二，这个栈槽承担什么角色。可能是：

\begin{itemize}
  \item 调试构建中为局部变量保留的位置。
  \item 寄存器不够时 spill 出去的临时值。
  \item 保存 callee-saved 寄存器。
  \item 为函数调用准备的栈上传参区域。
  \item 对齐填充或 red zone 中的临时空间。
  \item 编译器为了调试信息、异常展开或 sanitizer 插入的辅助数据。
\end{itemize}

例如：

\begin{lstlisting}
sub rsp, 16
mov dword ptr [rsp + 12], edi
mov eax, dword ptr [rsp + 12]
add rsp, 16
ret
\end{lstlisting}

在 \code{-O0} 中，这可能只是把参数 \code{x} 放到栈槽再读回来，方便调试器观察。优化后同样语义可能直接使用 \register{edi} 或 \register{eax}，栈槽完全消失。因此不能从 \code{-O0} 栈访问推断性能路径一定有内存访问。

栈访问还有 ABI 约束。函数调用前栈要满足对齐要求；\instruction{call} 会压入返回地址；\instruction{ret} 从栈顶取回返回地址；某些寄存器若属于 callee-saved，被调用者使用前要保存，返回前要恢复。这些内容将在 ABI 章节展开。本章先建立阅读规则：看到 \register{rsp} 或 \register{rbp} 的地址公式，不要马上写“局部变量”，而要结合函数序言、调用点、寄存器保存和优化级别判断。

\begin{warningbox}
栈槽是机器层存储位置，不是源码变量本身。一个源码变量可能没有栈槽，一个栈槽也可能在不同时间保存不同临时值。
\end{warningbox}

\section{\instruction{lea}：看起来像内存，实际不访问内存}

\instruction{lea} 是学习 x86 汇编必须尽早掌握的指令。它的名字是 load effective address，但它只计算地址表达式，不读取内存。

\begin{lstlisting}
lea rax, [rdi + rsi*4 + 8]
\end{lstlisting}

语义是：

\begin{lstlisting}
rax = rdi + rsi * 4 + 8
\end{lstlisting}

不是：

\begin{lstlisting}
rax = memory[rdi + rsi * 4 + 8]
\end{lstlisting}

对比：

\begin{lstlisting}
lea rax, [rdi + rsi*4 + 8]        ; compute address-like integer
mov eax, dword ptr [rdi + rsi*4 + 8] ; load 4 bytes from memory
\end{lstlisting}

编译器经常用 \instruction{lea} 做整数乘法常数：

\begin{lstlisting}
lea eax, [rdi + rdi*2]      ; eax = x * 3
lea eax, [rdi + rdi*4]      ; eax = x * 5
lea eax, [rdi*8 + 16]       ; eax = x * 8 + 16
\end{lstlisting}

这解释了为什么你在非指针代码里也会看到 \instruction{lea}：

\begin{lstlisting}[language=C++]
extern "C" int mul5_add1(int x)
{
    return x * 5 + 1;
}
\end{lstlisting}

可能生成：

\begin{lstlisting}
lea eax, [rdi + rdi*4 + 1]
ret
\end{lstlisting}

\begin{mentalmodel}
看到 \instruction{lea} 时先问：它的结果后面被当作地址使用，还是只是整数计算结果？\instruction{lea} 本身不产生 cache miss，因为它不访问内存。
\end{mentalmodel}

\instruction{lea} 的关键价值在于提醒我们：汇编语法中的方括号形式不总是意味着访存。对大多数指令，方括号表示内存操作数；但 \instruction{lea} 使用同样的寻址语法来计算有效地址，并把这个计算结果写进寄存器。它复用了硬件和编码中表达地址公式的能力，却不执行 load。

编译器喜欢 \instruction{lea}，因为它可以在一条指令中表达若干简单整数运算。它常用于计算指针，也常用于计算整数表达式，尤其是乘以 3、5、9 这类可以表示成基址加索引乘缩放的常数。读到 \instruction{lea} 时，不要被名字中的 load 误导；真正的判断标准是它是否读取内存。答案是否定的。

当然，\instruction{lea} 不是万能乘法器。它支持的缩放只有 1、2、4、8，并且表达式形状有限。复杂常数乘法可能需要多条 \instruction{lea}、移位、加减或 \instruction{imul}。这部分细节留到整数指令章节。第 6 章只要求你建立底线：\instruction{lea} 计算数值，不访问内存；它可能服务于地址，也可能服务于普通整数。

\section{flags：有些指令会留下条件信息}

x86 有 RFLAGS/EFLAGS。很多整数指令会更新 flags，例如 \instruction{add}、\instruction{sub}、\instruction{cmp}、\instruction{test}。后续条件跳转或条件移动会读取这些 flags。

本章先只建立直觉：

\begin{center}
\begin{tabular}{lll}
\toprule
flag & 名称 & 常见用途 \\
\midrule
ZF & zero flag & 结果是否为 0，相等比较 \\
SF & sign flag & 结果最高位是否为 1 \\
CF & carry flag & 无符号进位/借位 \\
OF & overflow flag & 有符号溢出 \\
\bottomrule
\end{tabular}
\end{center}

例如：

\begin{lstlisting}
cmp edi, esi
setg al
\end{lstlisting}

\instruction{cmp edi, esi} 本质上计算 \code{edi - esi} 来设置 flags，但不保存减法结果。\instruction{setg al} 根据 flags 把 \register{al} 设置为 0 或 1。控制流章节会系统讲条件码；这里只要记住：读汇编时，有些数据依赖不在通用寄存器里，而在 flags 里。

flags 是读汇编时最容易漏掉的隐式状态。很多指令没有把结果写进普通寄存器，却把条件信息写进 flags；后面的条件跳转、条件设置或条件移动再读取这些信息。若你只跟踪通用寄存器，就会看不懂为什么一条 \instruction{jle} 或 \instruction{setne} 能做出选择。

还要注意 flags 会被覆盖。两个会写 flags 的指令如果相邻出现，后者会覆盖前者留下的条件信息。于是条件跳转通常要紧跟产生条件的 \instruction{cmp}、\instruction{test} 或算术指令，除非中间指令不影响 flags。读控制流时，必须找到每个条件跳转读取的是哪一次 flags 写入。第 8 章会把这件事作为控制流恢复的核心方法。

\section{标签、函数和汇编器伪指令}

看一段 \code{clang++ -S} 生成的汇编，除了真正的机器指令，还会看到伪指令：

\begin{lstlisting}
    .text
    .globl  add2
    .type   add2,@function
add2:
    mov     eax, edi
    add     eax, esi
    ret
    .size   add2, .-add2
\end{lstlisting}

常见项：

\begin{center}
\begin{tabular}{p{0.25\linewidth}p{0.58\linewidth}}
\toprule
写法 & 含义 \\
\midrule
\code{.text} & 后面进入代码段 \\
\code{.globl add2} & 符号 \code{add2} 对链接器可见 \\
\code{.type add2,@function} & 标记符号类型是函数 \\
\code{add2:} & 标签，代表当前位置地址 \\
\code{.size add2, .-add2} & 记录函数大小 \\
\bottomrule
\end{tabular}
\end{center}

这些伪指令不是 CPU 执行的指令。它们给汇编器、链接器、调试器和异常展开机制使用。后续讲 ELF、链接、ABI 和 unwinding 时会展开。

初学阶段可以把汇编文件分成两类行：一类是 CPU 将来会执行的指令，一类是给工具链看的描述。标签介于两者之间：标签本身不是指令，但它命名了一个地址，跳转、调用、符号表和调试信息都可能引用它。伪指令通常以点开头，它们告诉汇编器如何组织输出，而不是告诉 CPU 做一次计算。

这个区分会影响反汇编阅读。你在 \code{.s} 文件中看到的伪指令，不一定会出现在 \code{objdump} 的指令列表里；反过来，反汇编会显示机器码地址和字节，这些在普通 \code{.s} 文件中未必直接出现。若报告要解释机器真正执行了什么，应以反汇编中的指令为准；若报告要解释符号、段和汇编器输入，则需要看 \code{.s}。

函数标签也不等于高级语言函数的全部语义。一个标签只是一个地址名字。函数边界、栈展开、调试行号、异常处理和符号可见性，需要伪指令、调试信息、ABI 和链接器共同配合。后续 ABI 章节会讲为什么函数调用不仅是 \instruction{call} 和 \instruction{ret}，还包括参数、返回值、栈对齐和寄存器保存约定。

\section{Intel 语法和 AT\&T 语法}

本书默认使用 Intel 语法，因为它和 Intel/AMD 手册、Compiler Explorer 的 Intel 输出、很多性能资料更接近。你仍然必须认识 AT\&T 语法，因为 GNU 工具链和很多历史资料会用它。

同一条指令：

\begin{lstlisting}
mov eax, dword ptr [rdi + rsi*4]      ; Intel
movl (%rdi,%rsi,4), %eax              ; AT&T
\end{lstlisting}

关键差异：

\begin{center}
\begin{tabular}{p{0.23\linewidth}p{0.31\linewidth}p{0.31\linewidth}}
\toprule
差异 & Intel & AT\&T \\
\midrule
操作数顺序 & destination, source & source, destination \\
寄存器写法 & \code{rax} & \code{\%rax} \\
立即数写法 & \code{42} & \code{\$42} \\
内存写法 & \code{[rdi + rsi*4]} & \code{(\%rdi,\%rsi,4)} \\
宽度表示 & \code{dword ptr} & 指令后缀 \code{movl} \\
\bottomrule
\end{tabular}
\end{center}

再看几个对应关系：

\begin{lstlisting}
Intel: mov rax, rdi
AT&T : movq %rdi, %rax

Intel: add eax, esi
AT&T : addl %esi, %eax

Intel: mov qword ptr [rsp + 8], rax
AT&T : movq %rax, 8(%rsp)
\end{lstlisting}

\begin{warningbox}
最容易错的是操作数顺序。Intel 是 \code{dst, src}，AT\&T 是 \code{src, dst}。读资料时先确认语法，否则会把数据流方向看反。
\end{warningbox}

\section{从函数签名建立入口状态}

读一个函数的汇编，第一步不是看第一条指令，而是看函数签名。函数签名告诉你参数个数、参数类型和返回类型；ABI 告诉你这些参数在函数入口时大致放在哪里。没有入口状态，后面的寄存器含义就没有根。

例如一个函数签名是“接收一个指针、一个整数下标、一个整数增量，返回整数”。在常见 System V AMD64 ABI 下，前三个整数或指针类参数通常从 \register{rdi}、\register{rsi}、\register{rdx} 进入；如果参数类型是 32 位 \code{int}，你常会看到它的低 32 位别名，例如 \register{esi} 或 \register{edx}。因此在读第一条指令前，你应该先写入口笔记：\register{rdi} 是基址指针，\register{esi} 是下标，\register{edx} 是增量，返回值最终应在 \register{eax}。

这个入口笔记不是最终答案，而是分析起点。编译器可以马上把参数复制到别的寄存器，可以把 32 位下标符号扩展到 64 位，可以把某个参数与常量合并，也可以完全删除某个未使用参数。只要你从入口状态开始逐条更新，就能跟上这些变化。反过来，如果你看到 \register{rax} 就猜“这一定是返回值”，很容易出错；\register{rax} 在函数中间也可能只是临时寄存器。

入口状态还要包含你不知道的内容。比如函数接收一个指针，你通常不知道它是否为空、是否对齐、指向多大对象、是否与另一个指针别名。汇编只告诉你程序按某个地址读取或写入，不能自动证明这个地址在 C++ 语义上合法。读汇编时，要把“寄存器里有一个地址值”和“这个地址一定指向有效对象”分开。后者需要语言规则、调用约定和程序上下文共同证明。

\begin{mentalmodel}
函数签名给出源码层入口，ABI 给出机器层入口，逐条指令更新当前状态。读汇编时不要从中间的某个寄存器名字开始猜，要从入口状态开始推。
\end{mentalmodel}

\section{确定事实、合理推论和待验证假设}

高质量汇编分析必须区分三类说法：确定事实、合理推论、待验证假设。三者都可以出现在报告里，但不能混在一起。

确定事实来自汇编文本和 ABI 规则。例如某条指令读取 \register{rdi}，写 \register{eax}；某个内存操作数的有效地址是 \register{rdi} 加 \register{rsi} 乘 4；某条 \instruction{mov} 访问 4 字节；函数返回值按 ABI 放在 \register{eax}。这些事实可以直接从指令、宽度标记和调用约定读出。

合理推论是在事实之上结合上下文得到的解释。例如 \code{[rdi + rsi*4]} 很可能是四字节元素数组访问；\code{[rdi + rax*8 + 16]} 可能是结构体数组中的字段；连续的 \instruction{cmp} 和条件跳转可能来自一个 \code{if}；\instruction{lea eax, [rdi + rdi*4 + 1]} 可能来自 \code{x * 5 + 1}。这些推论有依据，但不是机器文本本身直接写出来的。

待验证假设是还缺证据的解释。例如“这里慢是因为 cache miss”，“这里一定来自某个类成员函数”，“这个 \code{+ 24} 一定是第三个字段”，“这条间接调用一定调用某个虚函数”。这些说法可能正确，但需要更多证据：源码、类型布局、调试信息、符号、运行时采样、硬件计数器或调试器观察。没有证据时，应该把它写成假设，而不是结论。

这个区分看似麻烦，其实能显著提高分析质量。很多汇编报告的问题不是看错了指令，而是把推论写成事实。比如看到 \instruction{mov eax, dword ptr [rdi + 4]} 就说“读取结构体第二个字段”，这可能对，也可能不对。更严谨的写法是：这条指令从 \register{rdi} 加 4 的地址读取 4 字节；若 \register{rdi} 指向一个首字段大小为 4 字节且第二字段也是 4 字节的结构体，则它可能读取第二字段。

\begin{keyidea}
汇编分析的可信度来自证据等级。机器事实要写准，源码解释要标明推论条件，性能判断要说明还需要哪些测量证据。
\end{keyidea}

\subsection{从汇编反推源码不是唯一答案}

很多练习会让你“根据汇编写出可能的 C 代码”。这类练习有价值，但要知道：从汇编反推源码通常不是唯一答案。不同源码可以生成相同或相近汇编。

例如：

\begin{lstlisting}
lea eax, [rdi + rdi*4 + 1]
ret
\end{lstlisting}

它可能来自：

\begin{lstlisting}[language=C++]
return x * 5 + 1;
\end{lstlisting}

也可能来自：

\begin{lstlisting}[language=C++]
return x + 4 * x + 1;
\end{lstlisting}

还可能来自某个被内联函数、模板常量表达式、数组下标计算或优化器代数化简后的结果。汇编确定的是“返回 \register{edi} 乘 5 加 1 的 32 位结果”，不确定原始源码究竟写成哪种表达式。

再比如：

\begin{lstlisting}
mov eax, dword ptr [rdi + 16]
ret
\end{lstlisting}

它可能是结构体字段读取，也可能是数组偏移读取，也可能是某种对象头后面的值。只有当你知道函数签名、类型布局、调用上下文和源码，才能更有把握地解释。

因此，汇编阅读报告不应写得过度确定。可以使用下面的表达方式：

\begin{itemize}
  \item “机器事实是访问 \register{rdi} 加 16 处的 4 字节。”
  \item “若 \register{rdi} 指向 \code{S} 对象，且 \code{c} 字段偏移为 16，则这对应读取 \code{p->c}。”
  \item “当前尚未证明该地址一定来自结构体字段；需要源码或调试信息确认。”
\end{itemize}

这种谨慎不是多余文字，而是专业习惯。性能调试中，过度确定的错误解释比“不知道”更危险。

\subsection{证据强度排序}

同一个判断可能有不同强度的证据。比如判断某个函数是否在热路径中：

\begin{itemize}
  \item 源码中出现调用：证据弱，因为可能被内联、删除或走不到。
  \item 优化后反汇编有 \instruction{call}：证据较强，说明该二进制路径里有调用指令。
  \item profiling 显示采样落在该函数：证据更强，说明运行时确实执行并消耗时间。
  \item 控制输入和计数器后多次复现：证据最强，说明结论稳定。
\end{itemize}

判断一个地址公式来源也是类似：

\begin{itemize}
  \item 看到乘 4：只能说明地址中有 scale 4。
  \item 结合访问宽度 4 和函数签名 \code{int const*}：可以推论是 \code{int} 数组访问。
  \item 结合源码和调试信息：推论更强。
  \item 用 GDB 打印地址和内存，或用 \code{offsetof} 验证字段偏移：证据更完整。
\end{itemize}

本书要求你在报告里主动标注证据强度。尤其是性能解释，不要因为看见一条 load 就断言瓶颈是内存；不要因为看见一条 branch 就断言瓶颈是分支预测；不要因为指令条数变少就断言一定更快。汇编给出事实，瓶颈还需要测量。

\section{\code{-O0} 和 \code{-O3} 为什么差异巨大}

同一个 C++ 函数，在不同优化级别下可以完全不像同一个程序。看：

\begin{lstlisting}[language=C++]
extern "C" int f(int x)
{
    int y = x + 1;
    int z = y * 3;
    return z - 2;
}
\end{lstlisting}

\code{-O0} 的目标是容易调试，不是快。它可能生成栈帧，把局部变量放在栈上：

\begin{lstlisting}
push rbp
mov  rbp, rsp
mov  dword ptr [rbp - 4], edi
mov  eax, dword ptr [rbp - 4]
add  eax, 1
mov  dword ptr [rbp - 8], eax
mov  eax, dword ptr [rbp - 8]
imul eax, eax, 3
mov  dword ptr [rbp - 12], eax
mov  eax, dword ptr [rbp - 12]
sub  eax, 2
pop  rbp
ret
\end{lstlisting}

\code{-O3} 会做代数化简和寄存器分配，可能变成：

\begin{lstlisting}
lea eax, [rdi + rdi*2 + 1]
ret
\end{lstlisting}

因为：

\begin{lstlisting}
y = x + 1
z = y * 3 = 3*x + 3
return z - 2 = 3*x + 1
\end{lstlisting}

这个案例说明：

\begin{itemize}
  \item \code{-O0} 适合观察源代码结构、栈帧和调试信息。
  \item \code{-O3} 适合观察真实性能路径。
  \item 性能优化不能根据 \code{-O0} 汇编下结论。
  \item \code{-O3} 汇编可能已经和源码结构差很多，需要恢复语义而不是逐句对应。
\end{itemize}

\code{-O0} 和 \code{-O3} 的差异不是“一个粗糙、一个聪明”这么简单。它们服务的目标不同。\code{-O0} 尽量保留源代码级调试体验，让调试器能在某个源码行上看到局部变量，能单步经过看似对应的语句，能较少因为重排和删除让读者困惑。因此它常常把变量放入栈槽，频繁从内存读回，再写回内存。这个形态对学习栈帧和调试器很有用，但不能代表发布程序的热路径。

\code{-O3} 则以优化后的可观察行为为准，不承诺保留源码形状。局部变量可以消失，语句顺序可以改变，函数可以内联，分支可以合并，循环可以展开，甚至整个计算可以被常量折叠。读 \code{-O3} 汇编时，你的任务不是寻找“哪条指令对应哪行源码”，而是证明优化后的机器行为仍然实现了源码要求的结果，并找出真正执行的热路径。

两种输出都要看，但用途不同。\code{-O0} 适合观察编译器如何为调试构造保守形态，也适合理解栈帧和局部变量位置。\code{-O3} 适合观察性能相关路径，适合分析编译器如何利用寄存器、地址公式和代数化简。把两者混用会产生错误结论：用 \code{-O0} 判断性能会夸大内存访问，用 \code{-O3} 做源码级单步调试又会觉得变量“莫名其妙消失”。

\begin{warningbox}
报告中必须说明你分析的是哪种优化级别。没有优化级别，汇编分析不完整；用调试构建证明性能结论，通常不合格。
\end{warningbox}

\section{机器码字节和指令边界}

汇编文本是给人看的，机器最终执行的是字节。x86-64 是变长指令集，一条指令的长度不固定。短的指令可能只有 1 字节，带寄存器编码、位移、立即数、前缀或 REX 扩展的指令会更长。反汇编工具把字节流重新切分成指令，并显示地址、字节和助记符。

学习机器码字节不是为了让你手写编码表，而是为了建立两个关键直觉。第一，\register{rip} 顺序前进时，下一条指令地址等于当前地址加当前指令长度。控制流章节讲跳转时，这个事实会解释为什么相对跳转目标要根据当前指令结束位置计算。第二，指令长度会影响取指和前端压力。第二册深入 CPU 前端时，变长指令、解码带宽和代码尺寸会成为性能因素。

立即数和位移常常让指令变长。例如把一个小寄存器加到另一个寄存器，编码可能很短；如果指令还带有 32 位立即数，机器码就要额外存放这个常量；如果内存操作数带有较大的 displacement，也要把位移编码进指令。你在 \code{objdump} 里看到一条指令前面有很多十六进制字节，不要只看助记符，要问哪些字节来自操作码，哪些来自寄存器和寻址编码，哪些来自立即数或位移。

机器码字节还能帮助你理解反汇编和汇编文件的区别。汇编文件中的标签没有固定地址，直到汇编和链接后才有位置；反汇编中的地址和字节则来自已经生成的目标文件或可执行文件。若你要证明某个函数中真的存在某条机器指令，反汇编比源码和 \code{.s} 更接近证据。若你要解释编译器输出给汇编器的符号和伪指令，\code{.s} 又更直接。

\begin{deepdive}
本章不要求手工编码 x86 指令，但要求你不再把汇编文本当作最终执行物。机器码字节、指令长度和地址，是后续理解 \register{rip}、跳转、代码布局和前端性能的基础。
\end{deepdive}

\subsection{反汇编器怎样从字节流切出一条指令}

反汇编不是“把助记符翻回源码”，而是从某个起始地址开始，按 ISA 规则消费字节。对 x86-64，可以把反汇编器的最小状态机写成：

\begin{lstlisting}[numbers=none,caption={x86-64 反汇编的压缩状态机}]
StartAtAddress:
  pc = current byte address

ReadPrefixes:
  consume legacy prefixes, REX, VEX or EVEX if present

ReadOpcode:
  consume one or more opcode bytes

DecodeOperandEncoding:
  decide whether this opcode requires ModRM, SIB, disp or imm

ReadModRM:
  if required, split mod/reg/rm and maybe consume SIB

ReadDisplacementAndImmediate:
  consume disp8/disp32 and imm bytes according to opcode and ModRM

EmitInstruction:
  instruction_length = bytes_consumed
  next_pc = pc + instruction_length
\end{lstlisting}

这个流程有两个直接后果。第一，同一串字节如果从错误地址开始解码，后面的指令边界可能全错。第二，\register{rip} 相关的相对地址、短跳转和长跳转，都依赖 \code{next_pc}，也就是当前指令结束后的地址。反汇编器显示的每条指令地址不是装饰，它是计算控制流和 RIP-relative 地址的输入。

\begin{warningbox}
x86 没有天然的“每 4 字节一条指令”边界。反汇编必须从可信入口开始，例如符号表、函数入口、跳转目标或调试信息。把数据区、跳转表或字符串从中间拿来反汇编，工具也可能给出看似合法的指令文本，但那不代表程序真的会执行这些字节。
\end{warningbox}

\subsection{x86 编码的最小结构：prefix、opcode、ModRM、SIB、disp、imm}

x86-64 编码历史很长，完整规则足够写成手册。本书不要求背表，但必须知道一条指令的字节通常由哪些字段拼成。否则你无法解释为什么 \instruction{mov eax, edi} 是 2 字节，\instruction{mov rax, rdi} 是 3 字节，\instruction{mov eax, dword ptr [rdi + rsi*4 + 8]} 又更长。

一条常见指令可按下面顺序阅读：

\begin{lstlisting}[numbers=none,caption={x86-64 指令编码的常见字段顺序}]
legacy prefixes / REX prefix
opcode bytes
ModRM byte
SIB byte
displacement bytes
immediate bytes
\end{lstlisting}

\begin{center}
\small
\begin{tabular}{p{0.18\linewidth}p{0.31\linewidth}p{0.35\linewidth}}
\toprule
字段 & 作用 & 阅读重点 \\
\midrule
prefix & 改变宽度、段、重复、锁、向量扩展等 & \code{lock}、operand-size、REX、VEX/EVEX 都会影响语义或寄存器空间 \\
opcode & 指令种类和一部分方向/宽度信息 & 决定这是 \instruction{mov}、\instruction{add}、\instruction{ret} 还是别的指令族 \\
ModRM & 寄存器编号和寻址模式 & reg 字段、r/m 字段、mod 字段常决定寄存器或内存操作数 \\
SIB & scale-index-base 地址公式 & 用于表达 \code{base + index * scale + displacement} \\
displacement & 地址中的常量偏移 & 1/4 字节常见，RIP-relative 也用 disp32 \\
immediate & 指令中的字面常量 & 如 \code{add eax, 7} 中的 7，或 \code{mov eax, 0x1234} \\
\bottomrule
\end{tabular}
\end{center}

不是每条指令都有所有字段。\instruction{ret} 的 near return 是单字节 \code{c3}；\instruction{mov eax, edi} 有 opcode 和 ModRM；复杂内存寻址通常会有 ModRM、SIB 和 displacement；64 位寄存器或 \register{r8} 到 \register{r15} 往往需要 REX 前缀。

\begin{center}
\small
\begin{tabular}{p{0.25\linewidth}p{0.26\linewidth}p{0.33\linewidth}}
\toprule
字节示例 & 反汇编 & 字段切分 \\
\midrule
\code{c3} & \instruction{ret} & opcode only，长度 1 \\
\code{89 f8} & \instruction{mov eax, edi} & opcode \code{89} + ModRM \code{f8}，长度 2 \\
\code{48 89 f8} & \instruction{mov rax, rdi} & REX \code{48} + opcode \code{89} + ModRM \code{f8}，长度 3 \\
\code{05 78 56 34 12} & \instruction{add eax, 0x12345678} & opcode \code{05} + imm32，小端立即数，长度 5 \\
\bottomrule
\end{tabular}
\end{center}

这张表的作用不是让你背编码，而是训练“长度来自字段”的直觉。若某条指令有 5 字节，下一条指令地址就从当前地址加 5 开始；若某条 RIP-relative 访问带有 disp32，它的目标地址也要用指令结束地址计算。

\subsection{legacy prefix、REX、VEX、EVEX：前缀不是注释}

x86 前缀分几类。legacy prefix 继承自早期 x86，例如 \code{66} operand-size override、\code{f2}/\code{f3} repeat 或某些 SSE 标量浮点含义、segment override、\code{lock} 等。REX 是 x86-64 引入的单字节前缀，扩展寄存器和 64 位操作数宽度。VEX/EVEX 是 AVX/AVX-512 时代的多字节前缀，能编码更多寄存器、向量长度、三操作数形式、mask 和更多 opcode map 信息。

\begin{center}
\small
\begin{tabular}{p{0.18\linewidth}p{0.31\linewidth}p{0.35\linewidth}}
\toprule
前缀类别 & 常见字节/形式 & 读法 \\
\midrule
legacy & \code{66}、\code{f2}、\code{f3}、\code{f0} & 可能改变宽度、重复、锁语义或 SSE 标量指令含义 \\
REX & \code{40}--\code{4f} & W/R/X/B 扩展宽度和寄存器编号 \\
VEX & \code{c4 ...} 或 \code{c5 ...} & AVX 常见，携带向量长度、源寄存器、opcode map 等 \\
EVEX & \code{62 ...} & AVX-512 常见，携带 mask、广播、rounding、更多寄存器等 \\
\bottomrule
\end{tabular}
\end{center}

例如 \code{f3 0f 10 07} 常见反汇编为：

\begin{lstlisting}
movss xmm0, dword ptr [rdi]
\end{lstlisting}

这里的 \code{f3} 不能随便忽略。它不是普通字符串前缀，而是这条 SSE 指令编码的一部分，参与决定这是 \instruction{movss}，不是另一个 \code{0f 10} 指令形态。类似地，\code{c5 f4 58 c2} 的 \code{c5 f4} 是 VEX 前缀，决定 \instruction{vaddps ymm0, ymm1, ymm2} 的 AVX 三操作数和向量长度。

\subsection{REX 前缀：64 位宽度和扩展寄存器从哪里来}

x86-64 增加了 64 位寄存器和 \register{r8} 到 \register{r15}，但旧编码空间不够直接表达所有信息，于是引入 REX 前缀。REX 是 \code{0x40} 到 \code{0x4f} 的单字节前缀，低 4 位可理解为 W、R、X、B：

\begin{center}
\small
\begin{tabular}{p{0.12\linewidth}p{0.30\linewidth}p{0.42\linewidth}}
\toprule
位 & 作用 & 例子 \\
\midrule
W & 选择 64 位操作数宽度 & \code{48 89 f8} 常见为 \instruction{mov rax, rdi} \\
R & 扩展 ModRM.reg 字段 & 访问 \register{r8}--\register{r15} 的某些位置 \\
X & 扩展 SIB.index 字段 & 扩展 index 寄存器编号 \\
B & 扩展 ModRM.r/m 或 SIB.base 字段 & 扩展 base 或 r/m 寄存器编号 \\
\bottomrule
\end{tabular}
\end{center}

对比两个例子：

\begin{lstlisting}[numbers=none,caption={32 位和 64 位 move 的字节差异}]
89 f8        mov eax, edi
48 89 f8     mov rax, rdi
\end{lstlisting}

第二条多出的 \code{48} 是 REX.W=1，表示 64 位操作数宽度。没有它，\code{89 f8} 写的是 \register{eax}，并按 x86-64 规则清零 \register{rax} 高 32 位；有 \code{48}，写的是完整 \register{rax}。这不是反汇编显示风格差异，而是实际机器码语义差异。

扩展寄存器也依赖 REX。例如：

\begin{lstlisting}[numbers=none,caption={扩展寄存器需要 REX 扩展位}]
44 89 c0     mov eax, r8d
41 89 f8     mov r8d, edi
4c 89 c0     mov rax, r8
\end{lstlisting}

具体字节可能随汇编器选择的等价编码而变化，但原则稳定：看到 \register{r8}--\register{r15} 或 64 位通用寄存器操作时，要检查 REX 前缀。报告里只贴助记符不贴字节，就无法证明这些宽度和寄存器扩展细节。

\begin{deepdive}
REX 前缀还会影响传统高 8 位寄存器 \register{ah}、\register{bh}、\register{ch}、\register{dh} 的可编码性。现代 64 位代码通常避免主动使用这些寄存器。读到涉及 8 位寄存器和 REX 的手写汇编时，要格外谨慎；它不是纯语法问题，而是编码空间的历史限制。
\end{deepdive}

\subsection{ModRM：寄存器和寻址模式的三段位域}

很多 x86 指令在 opcode 后跟一个 ModRM 字节。ModRM 可拆成三段：

\begin{lstlisting}[numbers=none,caption={ModRM 字节的位域}]
bits 7..6: mod
bits 5..3: reg
bits 2..0: r/m
\end{lstlisting}

以 \code{89 f8} 为例，\code{89} 是 \instruction{mov r/m32, r32} 一类编码；\code{f8} 二进制是 \code{11111000}：

\begin{center}
\small
\begin{tabular}{p{0.18\linewidth}p{0.18\linewidth}p{0.46\linewidth}}
\toprule
字段 & 位值 & 含义 \\
\midrule
mod & \code{11} & 操作数是寄存器，不是内存寻址 \\
reg & \code{111} & 源寄存器编号 7，即 \register{edi} \\
r/m & \code{000} & 目标寄存器编号 0，即 \register{eax} \\
\bottomrule
\end{tabular}
\end{center}

所以 \code{89 f8} 可读成：把 reg 字段代表的 32 位寄存器 \register{edi} 移到 r/m 字段代表的 32 位寄存器 \register{eax}。同一个 ModRM 字段在不同 opcode 下可能扮演不同方向或含义，所以不能只看 ModRM，必须和 opcode 一起读。

\begin{center}
\small
\begin{tabular}{p{0.13\linewidth}p{0.35\linewidth}p{0.36\linewidth}}
\toprule
mod & r/m 含义概略 & 后续可能还有什么 \\
\midrule
\code{11} & r/m 是寄存器编号 & 通常没有 displacement \\
\code{00} & r/m 描述内存地址形态 & 特定 r/m 组合可能表示 SIB 或 RIP-relative/disp32 \\
\code{01} & 内存地址加 disp8 & 后面消费 1 字节 displacement \\
\code{10} & 内存地址加 disp32 & 后面消费 4 字节 displacement \\
\bottomrule
\end{tabular}
\end{center}

上表故意写成“概略”，因为完整 x86 编码有很多特例。但这已经足够建立阅读习惯：看到 ModRM 后，先看 \code{mod} 决定是寄存器还是内存；若是内存，再看 r/m 是否触发 SIB 或 RIP-relative；最后根据 \code{mod} 消费 displacement。

\subsection{SIB：base + index * scale + displacement 的来源}

当内存操作数需要 base/index/scale 形式时，ModRM 后可能跟一个 SIB 字节。SIB 拆成：

\begin{lstlisting}[numbers=none,caption={SIB 字节的位域}]
bits 7..6: scale
bits 5..3: index
bits 2..0: base
\end{lstlisting}

例如下面这类指令：

\begin{lstlisting}
mov eax, dword ptr [rdi + rsi*4 + 8]
\end{lstlisting}

常见编码形态会包含 opcode、ModRM、SIB 和 1 字节 displacement。你可以用：

\begin{lstlisting}[language=bash,caption={生成并拆一条 SIB 寻址指令}]
cat > /tmp/sib.s <<'ASM'
.intel_syntax noprefix
.text
.globl f
f:
    mov eax, dword ptr [rdi + rsi*4 + 8]
    ret
ASM
clang -c /tmp/sib.s -o /tmp/sib.o
objdump -drwC -Mintel /tmp/sib.o
\end{lstlisting}

报告里不要求你背出所有寄存器编号，但要能指出：某些字节在表达“这是内存操作数”；某个 SIB 字节在表达 scale=4、index=\register{rsi}、base=\register{rdi}；最后的 \code{08} 之类字节表达 displacement=8。这样你才能把汇编地址公式和机器码字节接上。

注意两个常见陷阱。第一，\register{rsp} 或 \register{r12} 作为 base 时，经常强制出现 SIB 字节，即使你没有显式 index。第二，某些 r/m 编码组合表示 RIP-relative 或 disp32 特例，不是普通 base 寄存器。x86 编码是历史兼容产物，不能用“看起来应该规则”替代反汇编证据。

\subsection{完整手拆一条：48 8b 44 b7 08}

下面这条来自最小 x86-64 目标文件：

\begin{lstlisting}[numbers=none]
0: 48 8b 44 b7 08    mov rax, qword ptr [rdi + 4*rsi + 0x8]
\end{lstlisting}

按字段拆：

\begin{center}
\small
\begin{tabular}{p{0.15\linewidth}p{0.20\linewidth}p{0.49\linewidth}}
\toprule
字节 & 字段 & 含义 \\
\midrule
\code{48} & REX & W=1，选择 64 位操作数宽度 \\
\code{8b} & opcode & \instruction{mov r64, r/m64} 形态，方向是从 r/m 读到 reg \\
\code{44} & ModRM & mod=\code{01}，reg=\code{000}，r/m=\code{100}；目标 reg 是 \register{rax}，r/m=\code{100} 表示后面有 SIB，且有 disp8 \\
\code{b7} & SIB & scale=\code{10} 表示乘 4，index=\code{110} 是 \register{rsi}，base=\code{111} 是 \register{rdi} \\
\code{08} & disp8 & 地址再加 8 \\
\bottomrule
\end{tabular}
\end{center}

所以它不是凭文本猜出来的 \code{[rdi + rsi*4 + 8]}，而是字节字段共同证明的地址公式：

\begin{lstlisting}[numbers=none]
effective_address = rdi + rsi * 4 + 8
load_width = 8 bytes
destination = rax
instruction_length = 5 bytes
next_rip = current_rip + 5
\end{lstlisting}

注意最后两行。读反汇编时，指令长度和下一条指令地址同样是证据。它们会进入 RIP-relative 地址、跳转目标、patch 点、断点覆盖和前端取指分析。

\subsection{RIP-relative：位移相对下一条指令，不是当前地址}

x86-64 常用 RIP-relative 访问全局数据、常量或 GOT 表项。语法可能显示为：

\begin{lstlisting}
mov eax, dword ptr [rip + 0x2f3a]
\end{lstlisting}

这里的有效地址不是“当前指令起始地址 + 0x2f3a”，而是“下一条指令地址 + 0x2f3a”。原因是 x86 相对位移通常以当前指令结束位置作为基准。跳转、调用和 RIP-relative data access 都要保留这个直觉。

\begin{lstlisting}[numbers=none,caption={RIP-relative 地址计算模板}]
instruction_start = A
instruction_length = L
next_rip = A + L
encoded_disp32 = D
effective_address = next_rip + sign_extend_64(D)
\end{lstlisting}

这会直接影响调试。若 \register{rip} 附近反汇编显示访问 \code{[rip + disp]}，你要先看该指令长度，再算目标地址；如果目标文件尚未链接，还要看重定位记录，因为 disp 可能只是占位。正确报告应同时贴：

\begin{lstlisting}[numbers=none,caption={RIP-relative 证据包}]
objdump -drwC -Mintel file.o
readelf -r file.o
nm -n file.o
\end{lstlisting}

只贴 \code{mov eax, [rip+0x0]} 往往是不完整的。目标文件里 \code{0x0} 可能只是等待链接器修补的 relocation 占位；链接后可执行文件里才会出现最终位移。

\subsection{位移和立即数：小端序也会出现在指令里}

x86 是小端机器，不只数据内存如此，指令中的多字节 displacement 和 immediate 也按小端顺序编码。例如：

\begin{lstlisting}
add eax, 0x12345678
\end{lstlisting}

机器码里的立即数字节通常会按 \code{78 56 34 12} 出现，而不是按文本顺序 \code{12 34 56 78}。这能解释为什么反汇编里的十六进制字节和助记符常量看起来“反过来”。

\begin{lstlisting}[language=bash,caption={观察立即数小端编码}]
cat > /tmp/imm.s <<'ASM'
.intel_syntax noprefix
.text
.globl g
g:
    add eax, 0x12345678
    ret
ASM
clang -c /tmp/imm.s -o /tmp/imm.o
objdump -drwC -Mintel /tmp/imm.o
\end{lstlisting}

读机器码字节时要把“字节流顺序”和“数值显示”分开。CPU 按地址顺序取字节；解码器按 ISA 规则把后续若干字节解释成 little-endian 的位移或立即数；反汇编器再把数值显示成人类习惯的十六进制。

\subsection{AVX/VEX 字节长什么样：先识别，不急着手拆}

现代 SIMD 指令常用 VEX 或 EVEX 前缀，字段比普通 REX/ModRM 复杂得多。第一册不要求手工拆完 VEX 每一位，但必须能识别“这些前缀字节决定了 AVX 语义”。例如：

\begin{lstlisting}[numbers=none]
c5 f4 58 c2       vaddps ymm0, ymm1, ymm2
c4 e2 75 b8 c2    vfmadd231ps ymm0, ymm1, ymm2
\end{lstlisting}

第一条以 \code{c5} 开始，是 2 字节 VEX 前缀；第二条以 \code{c4} 开始，是 3 字节 VEX 前缀。它们把传统 SSE 的二操作数模型扩展成 AVX 的三操作数模型，并编码向量长度、源寄存器和 opcode map 等信息。报告里如果声称“这是 AVX 指令”“这是三操作数 packed float 加法”，至少要贴出反汇编文本和机器码字节，说明 VEX 前缀存在。

到第二册做 SIMD 性能时，再继续追 VEX/EVEX 细节：向量长度位、mask、broadcast、rounding、zeroing/merging、EVEX 压缩 displacement 等。第一册现在只建立边界：看到 \code{c4}、\code{c5}、\code{62} 这类前缀，不要把它们当成可忽略的“乱码”；它们正是现代向量指令语义的一部分。

\subsection{编码证据报告：最少要贴什么}

如果一份报告声称“这里用了 RIP-relative”“这里是 64 位寄存器写”“这里访问 \register{r8}”“这里的分支目标是某个地址”，最少应给出：

\begin{itemize}
  \item 反汇编命令和工具版本。
  \item 目标架构和文件格式，例如 \code{elf64-x86-64}。交叉编译或在非 x86-64 主机上分析 x86 文件时尤其要写。
  \item 指令地址、机器码字节和反汇编文本。
  \item 若有符号或全局访问，附 \code{readelf -r} 的重定位记录。
  \item 对 REX、ModRM、SIB、disp、imm 中和结论相关的字段做最小解释。
  \item 若是 SSE/AVX/AVX-512，说明 legacy/VEX/EVEX 编码、向量宽度和操作数数量。
  \item 写清哪些是机器码直接证明的事实，哪些是结合源码或符号表的推论。
\end{itemize}

例如：

\begin{lstlisting}[numbers=none,caption={编码证据记录示例}]
bytes:
  48 8b 44 b7 08
text:
  mov rax, qword ptr [rdi + rsi*4 + 8]
claim:
  REX.W proves 64-bit destination width.
  SIB proves scaled-index addressing is present.
  objdump text gives the authoritative decoded address formula.
boundary:
  This proves the machine address formula, not the C++ source type.
\end{lstlisting}

最后一句很重要。编码可以证明“机器码如何寻址”，不能单独证明“这是某个结构体字段”或“访问一定安全”。类型、安全性和性能瓶颈仍要回到源码、对象布局、运行时输入和测量证据。

\section{读汇编的固定流程}

以后每次读一个函数，按下面流程做笔记：

\begin{enumerate}
  \item 写出函数签名和 ABI 参数位置。
  \item 标出每个入口寄存器的含义，例如 \register{rdi} 是第一个指针，\register{esi} 是第二个 \code{int}。
  \item 从上到下维护寄存器含义表。
  \item 对每条内存操作写出有效地址公式。
  \item 标出访问宽度和读写字节范围。
  \item 标出哪些指令更新 flags，后面哪里读取 flags。
  \item 区分真正 load/store 和纯地址计算 \instruction{lea}。
  \item 最后再尝试恢复 C-like 伪代码。
\end{enumerate}

例如：

\begin{lstlisting}
movsxd rsi, esi
mov    eax, dword ptr [rdi + rsi*4 + 8]
add    eax, 1
ret
\end{lstlisting}

笔记应该像这样：

\begin{lstlisting}
entry:
    rdi = base pointer p
    esi = int index i

movsxd rsi, esi
    rsi = sign_extend_64(i)

mov eax, dword ptr [rdi + rsi*4 + 8]
    address = p + i * 4 + 8
    load width = 4 bytes
    possible source = p[i + 2] if p is int*

add eax, 1
    eax = loaded_value + 1

ret
    return eax
\end{lstlisting}

这种笔记比“这大概是数组访问”更有价值。高性能工程里，证据链必须精确到地址公式和访问宽度。

这个流程看起来繁琐，但它的作用是减少猜测。初学者常常一看到熟悉的指令就跳到结论，例如看到 \instruction{lea} 就说“取地址”，看到 \code{+ 8} 就说“第三个元素”，看到 \instruction{ret} 就说“返回上一行”。固定流程强迫你先写事实：入口寄存器是什么，当前寄存器含义如何变化，内存地址公式是什么，访问宽度是多少，flags 在哪里产生和消费。事实写清楚后，结论通常自然出现。

流程中的最后一步才是恢复 C-like 伪代码。伪代码应该尽量贴近机器行为，而不是强行恢复原始源码格式。比如一个结构体数组字段更新，伪代码可以写成“读取地址 \code{base + i * 12 + 8} 的 4 字节整数，加上 \code{delta}，写回同一地址，再把结果作为返回值”。这比直接写“\code{items[i].count += delta}”更能展示你真的理解了地址公式和访存行为。

如果某一步缺证据，就把缺口写出来。例如你不知道 \register{rdi} 的源类型，可以先说它是第一个参数指针；如果没有源码和调试信息，不能断言它一定是某个结构体。工程报告不要求每个推论都完美，但要求推论边界清楚。

\section{案例：从 C++ 到汇编逐条解释}

创建函数：

\begin{lstlisting}[language=C++]
struct Item {
    int id;
    float score;
    int count;
};

extern "C" int update(Item* items, int i, int delta)
{
    items[i].count += delta;
    return items[i].count;
}
\end{lstlisting}

常见布局：

\begin{center}
\begin{tabular}{llll}
\toprule
字段 & offset & 大小 & 说明 \\
\midrule
\code{id} & 0 & 4 & \code{int} \\
\code{score} & 4 & 4 & \code{float} \\
\code{count} & 8 & 4 & \code{int} \\
\bottomrule
\end{tabular}
\end{center}

\code{sizeof(Item)} 通常是 12。因为 x86 scale 没有 12，编译器可能使用 \instruction{lea} 生成 \code{i * 3}，再乘 4：

\begin{lstlisting}
movsxd rax, esi
lea    rax, [rax + rax*2]
add    dword ptr [rdi + rax*4 + 8], edx
mov    eax, dword ptr [rdi + rax*4 + 8]
ret
\end{lstlisting}

逐条解释：

\begin{lstlisting}
entry:
    rdi = items
    esi = i
    edx = delta

movsxd rax, esi
    rax = sign_extend_64(i)

lea rax, [rax + rax*2]
    rax = i * 3

add dword ptr [rdi + rax*4 + 8], edx
    address = items + (i * 3) * 4 + 8
            = items + i * 12 + 8
            = addr(items[i].count)
    load old 4-byte count
    add delta
    store new 4-byte count
    update flags

mov eax, dword ptr [rdi + rax*4 + 8]
    reload 4-byte count into eax as return value

ret
\end{lstlisting}

为什么第二条 \instruction{mov} 要 reload？因为 \instruction{add dword ptr [mem], edx} 是 read-modify-write 内存指令，它把结果写回内存，但不把结果写入 \register{eax}。返回值需要在 \register{eax}，所以还要再 load。编译器也可能生成不同形式，比如先 load 到寄存器、加、store、再 return 寄存器。你要解释自己的输出。

\begin{deepdive}
read-modify-write 指令表面是一条汇编，但微架构内部通常涉及 load、ALU、store，并会占用 load/store 相关资源。后面讲 uops、ports、store buffer 和原子指令时，这一点会变得非常重要。
\end{deepdive}

\section{汇编阅读报告的最低质量线}

本章实验要求写报告，不是为了形式，而是训练证据链。一份合格的汇编阅读报告至少要让另一个读者在没有你口头解释的情况下复现你的判断。

报告开头应写清楚源码文件、编译器版本、编译命令、优化级别、目标平台和使用的反汇编工具。然后给出函数签名和 ABI 入口状态。对每个被分析函数，先列出参数寄存器和返回寄存器，再摘出核心指令。不要把整个反汇编大段粘进去；应该选择与分析相关的指令，并说明省略了什么。

逐条注释要包含读写状态。最低要求是：操作数类型、读取寄存器或内存、写入寄存器或内存、访问宽度、是否更新 flags、指令后的寄存器含义变化。遇到内存操作数，必须写有效地址公式；遇到结构体或数组解释，必须说明元素大小或字段偏移依据；遇到 \instruction{lea}，必须明确它是否访问内存；遇到 \instruction{cmp} 或 \instruction{test}，必须说明后续哪条指令读取 flags。

报告还要包含推论边界。能从汇编直接看出的，写成事实；需要源码类型支持的，写成“若某类型布局成立，则”；需要运行时数据支持的，写成待验证假设。比如“这条指令从 \register{rdi}+16 读取 4 字节”是事实；“这可能是读取结构体字段 \code{c}”是推论；“这里慢是因为缓存未命中”则需要性能计数器或实验支持。

最后，报告要说明 \code{-O0} 和 \code{-O3} 的用途差异。如果你同时展示两者，应分别解释它们服务什么问题：调试构建帮助理解栈帧和变量可见性，优化构建帮助理解真实性能路径。不能把两份汇编混在一起得出一个没有上下文的性能判断。

\begin{keyidea}
汇编阅读报告不是“贴代码加中文翻译”。它应当证明你能从函数签名和 ABI 建立入口状态，从指令读写恢复数据流，从内存操作数恢复地址公式，并清楚区分事实、推论和假设。
\end{keyidea}

\section{把汇编读成证据链，而不是读成故事}

真正可靠的汇编阅读，最后应该形成一条证据链。证据链不是把源码、汇编、反汇编、运行输出都贴上去，而是说明这些材料分别证明了什么。源码证明程序员写下了哪些语言层意图；编译命令证明编译器在什么规则下工作；\filepath{.s} 文件证明编译器交给汇编器的文本是什么；目标文件证明哪些机器码、符号和重定位已经生成；链接后的可执行文件证明链接器最终把符号绑定到哪里；调试器或运行日志证明某次运行时寄存器和内存状态怎样变化。每一层都有自己的边界，混成一团就会出现看似合理、实则无法复现的结论。

例如，你在 \filepath{.s} 文件里看到：

\begin{lstlisting}
call external_add@PLT
\end{lstlisting}

这可以证明编译器输出中存在一次对外部符号的调用意图，但不能直接证明运行时一定跳到哪个共享库地址。要证明目标文件阶段还没有确定最终地址，需要看 \code{objdump -drwC} 的重定位旁注或 \code{readelf -r} 的重定位项。要证明链接后地址，需要看可执行文件反汇编或动态链接信息。要证明某次运行绑定到哪个共享库实现，还要看 \code{ldd}、\code{LD_DEBUG}、\code{gdb} 中的符号解析，或者运行时映射。把这些层次分开，报告才有审查价值。

再看一个更常见的例子。你在反汇编里看到：

\begin{lstlisting}
mov eax, dword ptr [rdi + rsi*4 + 8]
\end{lstlisting}

这条指令直接证明的事实只有几个：读取 \register{rdi}、\register{rsi}，计算地址 \code{rdi + rsi*4 + 8}，从该地址读取 4 字节，写入 \register{eax}，并因为写 32 位寄存器而清零 \register{rax} 高 32 位。它不能单独证明这里一定是 \code{p[i + 2]}，也不能单独证明 \register{rdi} 一定指向 \code{int} 数组。若你有函数签名 \code{int f(int const* p, long i)}，则可以把 \register{rdi} 和 \register{rsi} 分别解释成 \code{p} 和 \code{i}，并推导 \code{+ 8} 可能是两个 \code{int} 元素的偏移。若源码是结构体数组，则 \code{+ 8} 也可能是字段 offset。结论要依赖签名、类型布局和上下文。

\begin{mentalmodel}
证据链写法的核心句式是：“从某个材料可以直接看到什么；结合哪个规则可以推出什么；仍然不能断言什么”。这比直接写“这一行表示访问数组”更慢，但更接近工程分析。
\end{mentalmodel}

\subsection{四层材料的常见用途}

第一层材料是编译输入和命令。它包括源码、头文件、编译器版本、优化级别、目标 CPU、是否使用 \code{-fPIC}、是否开启调试信息、是否保留帧指针、是否链接时优化。没有这一层，汇编输出没有环境。相同源码在 \code{-O0}、\code{-O2}、\code{-O3} 下可以完全不同；相同 \code{-O3} 在不同编译器版本下也可能选择不同指令形态。报告里不要只写“生成汇编”，而要写完整命令。

第二层材料是 \filepath{.s}。它适合观察编译器给汇编器的结构意图：函数标签、section 切换、伪指令、对齐、局部标签、常量池、展开信息和符号可见性。对于理解“编译器为什么把某个常量放进 \code{.rodata}”“为什么某个函数被标为全局符号”“为什么循环入口前有对齐伪指令”，\filepath{.s} 往往比反汇编更直观。

第三层材料是目标文件。它适合观察已经编码的机器码、符号表、重定位表和 section header。目标文件仍未完成所有地址绑定。外部符号、位置无关代码中的全局对象访问、跨 section 引用，都可能留下重定位项。读目标文件时要习惯同时看三类输出：反汇编、符号表、重定位表。只看反汇编会漏掉“这条 \instruction{call} 的目标还要链接器修补”这种关键信息。

第四层材料是链接后产物和运行时观察。链接后，函数和对象有了更具体的地址布局，PLT/GOT、动态重定位、可执行文件入口、只读段和可写段关系更清楚。运行时，地址还会受到 ASLR、动态加载器、共享库版本和内核映射影响。若分析目标是性能或崩溃，最终还需要运行时证据：寄存器、栈、内存、分支路径、采样位置、性能计数器。反汇编解释的是可能执行什么，运行时证据说明某次或某类输入实际执行了什么。

\subsection{不要把不同层次的地址混用}

地址是初学者最容易混淆的对象。汇编文件里的标签没有最终虚拟地址；目标文件里的 section 偏移不是进程运行时地址；链接后可执行文件中的虚拟地址可能还会被 PIE 和 ASLR 改变；GDB 里看到的地址是某次运行的加载地址。它们都可以写成十六进制数，但含义不同。

一个严谨的报告应该说明地址属于哪一层。例如“目标文件 \code{foo.o} 的 \code{.text} section 中 offset \code{0x20} 处有一条 \instruction{call}”和“运行时进程地址 \code{0x7f...} 处执行 \instruction{call}”不是同一句话。前者用于解释目标文件结构，后者用于解释某次运行。若要把两者对应起来，需要知道 section 加载基址、符号地址和调试信息。

同样，跳转目标也有层次。目标文件反汇编里，某个 \instruction{jmp} 可能显示为相对当前 section 的标签；链接后可能显示为函数内偏移；运行时采样工具可能显示为模块名加偏移。不要因为三个输出显示的十六进制不同，就认为程序变了；也不要因为助记符相同，就认为地址已经一致。高质量分析要说明“这些地址如何对应”。

\section{从机器状态表开始读陌生函数}

陌生函数最稳的读法不是先猜源码，而是先建立机器状态表。机器状态表记录每个关键时刻寄存器、内存别名、flags 和栈指针的大致含义。它不需要记录所有寄存器，但必须记录当前函数会读取、写入或跨调用保持的状态。

函数入口状态来自 ABI。对于常见 System V AMD64 整数参数，\register{rdi}、\register{rsi}、\register{rdx}、\register{rcx}、\register{r8}、\register{r9} 是前六个整数或指针参数的通道；返回整数通常在 \register{rax}。如果函数使用浮点参数，还要看 \register{xmm0} 等寄存器。入口状态不是从第一条 \instruction{mov} 猜出来的，而是由函数签名和 ABI 共同给出。

接着从上到下维护“当前含义”。例如：

\begin{lstlisting}
entry:
    rdi = p
    esi = i
    edx = delta

movsxd rax, esi
    rax = sign_extend_64(i)

lea rcx, [rax + rax*2]
    rcx = i * 3

mov eax, dword ptr [rdi + rcx*4 + 8]
    eax = load32(p + i * 12 + 8)
\end{lstlisting}

这种记录的价值在于，它把“寄存器当前是什么”从脑子里拿出来，变成可以审查的材料。只要中间某一步写错，后续地址公式就会暴露矛盾。比如你把 \register{rcx} 误记成 \code{i * 4}，下一条地址就会变成 \code{p + i * 16 + 8}，与结构体大小不符。机器状态表能迫使你及时发现这种错误。

\subsection{状态表要记录失效，而不只记录赋值}

寄存器含义不是永久有效的。任何写寄存器的指令都会覆盖旧含义；任何普通函数调用都可能破坏 caller-saved 寄存器；任何写 flags 的指令都会覆盖前一个比较结果；任何 store 都可能改变某个内存位置。状态表必须记录这些失效。

例如：

\begin{lstlisting}
cmp edi, esi
lea eax, [rdi + 1]
jl .Lless
\end{lstlisting}

这里 \instruction{lea} 不写 flags，所以 \instruction{jl} 仍然读取 \instruction{cmp} 产生的 flags。若中间换成：

\begin{lstlisting}
cmp edi, esi
add eax, 1
jl .Lless
\end{lstlisting}

那么 \instruction{add} 会写 flags，\instruction{jl} 读取的就不再是 \instruction{cmp} 的结果。真实编译器不会无意生成这种错误，但手写汇编和人工改写很容易犯。状态表若只记录寄存器值，不记录 flags 生命周期，就读不出这个问题。

函数调用也是同样道理。看到 \instruction{call} 之后，除非 ABI 或内联上下文能证明，否则不要继续相信 \register{rax}、\register{rcx}、\register{rdx}、\register{rdi}、\register{rsi}、\register{r8} 到 \register{r11} 的旧值还在。若调用前某个值调用后还要用，编译器通常会把它放到 callee-saved 寄存器、栈槽，或者在调用后重新计算。报告中若跨 \instruction{call} 使用旧寄存器含义，必须解释保存证据。

\subsection{内存状态要写别名假设}

寄存器状态比较容易记录，内存状态更难，因为多个地址可能指向同一对象。假设你看到：

\begin{lstlisting}
mov dword ptr [rdi], 1
mov eax, dword ptr [rsi]
\end{lstlisting}

如果 \register{rdi} 和 \register{rsi} 可能相等，那么第二条 load 可能读到刚写入的 1；如果根据函数契约、类型或 restrict-like 信息可以证明不别名，第二条 load 才能独立解释。汇编本身只显示两个地址来源，不显示语言层别名证明。报告要写清楚：这里是否假设两个指针不重叠，证据是什么。

别名问题在结构体字段、数组切片和手写汇编包装函数中尤其常见。两个地址公式看起来不同，不代表一定不重叠；同一个 base 加不同 offset，若访问宽度交叉，也可能部分重叠。读汇编时，内存记录至少要包含 base、index、scale、displacement、访问宽度和已知对象范围。只有这样，才能判断 store 是否可能影响后续 load。

\begin{warningbox}
“寄存器含义表”如果不记录 flags 失效、调用破坏和内存别名，只是半张表。完整的机器状态分析必须同时覆盖寄存器、flags、栈和内存。
\end{warningbox}

\section{用反例校正常见直觉}

汇编入门最危险的不是不会，而是会了一半后形成过度简化的直觉。本节列出几个经常导致错误报告的反例。

第一，\instruction{lea} 不等于“取地址”。它的操作数写法像内存，但不访问内存。编译器常用它做加法、常数乘法和地址中间量。若 \instruction{lea eax, [rdi + rsi*4]} 出现在整数函数里，可能只是计算 \code{x + y * 4}，并没有形成可解引用指针。是否是地址，取决于后续如何使用结果，而不取决于方括号外形。

第二，\instruction{mov} 不等于 C++ 赋值。C++ 赋值可能调用重载赋值运算符，可能写多个字段，可能涉及对象生命周期；\instruction{mov} 只是复制固定位宽的位模式。反过来，一个 C++ 赋值也可能完全消失，因为值被常量传播或保存在寄存器里。把 \instruction{mov} 机械翻译成“赋值语句”，会漏掉语言和机器之间的巨大差异。

第三，一个栈槽不等于一个源码局部变量。在 \code{-O0} 下，编译器可能为调试可见性给每个局部变量安排栈槽；在 \code{-O3} 下，多个临时值可能复用同一栈位置，或者根本没有栈槽。栈槽是编译器实现选择，不是 C++ 对象存在的唯一证据。对象生命周期要结合源码语义和调试信息判断，不能只看 \code{[rbp - 4]}。

第四，反汇编的函数边界不总是源代码函数边界。内联会让被调用函数的指令进入调用者；尾调用会让一个函数以 \instruction{jmp} 结束；链接器可能插入 PLT、thunk 或安全检查序列；异常展开和冷路径分离会让函数代码不完全连续。看到一个符号名，只能说明工具按某种符号信息展示了一个范围，不能保证它和源码函数体一一对应。

第五，运行输出正确不代表汇编接口正确。一个手写汇编函数可能碰巧返回正确 \register{eax}，但破坏了 \register{rbx}；一个函数可能在小输入下栈未对齐也没崩溃，但在调用库函数或开启向量化后出错；一个参数扩展错误可能在正数测试中看不出来，一遇到负数或高位非零就失败。汇编正确性要检查契约，不只看样例输出。

\section{汇编注释的层级写法}

汇编注释应该分层，而不是每行后面写一句模糊中文。建议使用三层注释。

第一层是机器语义注释。它只描述指令读写状态。例如“从 \code{[rdi + rsi*4]} 读取 4 字节到 \register{eax}”“写 \register{eax} 并清零 \register{rax} 高 32 位”“更新 ZF、SF、OF、CF 等 flags”。这一层要尽量客观，不引入源码猜测。

第二层是类型和布局注释。它结合函数签名、结构体布局、数组元素大小，把地址公式解释成对象访问。例如“若 \register{rdi} 是 \code{int const* p}，\register{rsi} 是下标 \code{i}，则访问 \code{p[i]}”“若结构体大小为 12，offset 8 是字段 \code{count}，则访问 \code{items[i].count}”。这一层要说明前提。

第三层是意图和性能注释。它解释编译器为什么可能选择这种形态，以及这个形态对后续分析有什么影响。例如“使用 \instruction{lea} 是为了计算 \code{i * 3}，不写 flags”“内存目的 \instruction{add} 是 RMW，结果没有自动进入返回寄存器”“这里选择 \instruction{movzx} 是因为无符号字节提升到 \code{int}”。这一层可以提出推论，但要避免把推论写成事实。

一个高质量注释块可以写成：

\begin{lstlisting}
movzx eax, byte ptr [rdi + rsi]
    machine:
        load 1 byte from rdi + rsi
        zero-extend to 32 bits
        write eax, clear high half of rax
    layout:
        if rdi = p and rsi = i, this reads p[i]
    source-level:
        matches uint8_t load promoted to int
    boundary:
        does not prove i is in bounds; bounds are a C++ precondition
\end{lstlisting}

这种写法比单行“读取数组元素”长，但它一次性说明了机器宽度、扩展方式、类型依据和边界条件。对复杂函数，不需要每条指令都写四层，但关键 load、store、flags、call、ret、地址计算和跨边界指令必须达到这个精度。

\section{优化汇编阅读：先接受源码形状已经消失}

读 \code{-O3} 汇编时，最大的心理障碍是源码形状消失。局部变量消失、函数消失、循环形态改变、分支变成条件移动、乘法变成 \instruction{lea}、小函数被内联、常量被提前算出，这些都不是异常，而是优化构建的正常结果。优化器的目标不是保留源码外形，而是在语言规则允许的范围内生成更合适的机器程序。

因此，优化汇编阅读要从“机器现在做什么”开始，而不是从“源码原来长什么样”开始。你可以把分析分成三轮。

第一轮只读机器事实。列出函数边界、入口寄存器、核心基本块、load/store、call、ret、flags 生产和消费。此时不要急着写 C++ 伪代码。机器事实应能在没有源码的情况下成立。

第二轮结合源码类型。把入口寄存器映射到参数，把 displacement 映射到字段 offset，把 scale 映射到元素大小，把条件码映射到 signed/unsigned 语义，把返回寄存器映射到返回类型。此时可以写“若源码类型为某结构，则该地址公式对应某字段”。

第三轮解释优化形态。问编译器为什么可以这样做：是否内联、是否常量传播、是否删除无用计算、是否利用未定义行为前提、是否因为别名信息受限而保留 load、是否因为 ABI 边界无法继续优化。优化解释必须建立在前两轮事实之上。

\subsection{案例：函数消失不等于代码没执行}

考虑：

\begin{lstlisting}[language=C++]
static int twice(int x)
{
    return x * 2;
}

extern "C" int caller(int x)
{
    return twice(x) + 1;
}
\end{lstlisting}

优化后你可能找不到独立的 \code{twice} 符号，\code{caller} 中只剩：

\begin{lstlisting}
lea eax, [rdi + rdi + 1]
ret
\end{lstlisting}

这不表示 \code{twice} “没有被执行”，而是它的语义被内联并合并到调用者中。报告应写：“目标二进制中没有独立 \code{twice} 函数符号；\code{caller} 的指令直接计算 \code{x * 2 + 1}；因此当前构建下 \code{twice} 的函数调用开销不存在。”不要写“编译器把函数删错了”。

如果实验目标是测函数调用开销，就需要阻止内联或把函数放到单独编译单元，并说明这样做改变了实验问题。如果实验目标是真实发布性能，则内联是正常路径，应接受它并分析内联后的机器码。

\subsection{案例：常量折叠让循环消失}

考虑：

\begin{lstlisting}[language=C++]
extern "C" int sum_const()
{
    int s = 0;
    for (int i = 0; i < 100; ++i) {
        s += i;
    }
    return s;
}
\end{lstlisting}

优化后可能直接返回常量：

\begin{lstlisting}
mov eax, 4950
ret
\end{lstlisting}

机器事实是：当前二进制没有循环，只有一次立即数写入返回寄存器。源码解释是：循环边界和每次加数都是编译期常量，且无副作用，编译器可以在编译期算出结果。性能解释是：这个函数不能用于测循环加法吞吐，因为被测循环不存在。

这类例子是 benchmark 中最常见的陷阱。你以为测了一亿次加法，其实测的是返回一个常量。防止它的方法不是关闭所有优化，而是让输入来自运行时、让结果被观察、并用反汇编确认热路径存在。关闭优化会把实验变成调试构建性能，也不是你想测的发布路径。

\subsection{案例：符号存在不代表热路径调用它}

有时 \code{nm} 显示某个函数符号存在，但调用点仍然没有调用它。原因可能是其他调用点需要该符号，或函数被导出供外部使用，但当前热路径已经内联了它。符号表只能证明“二进制中有这个符号”，不能证明“当前输入执行了这个符号”。

要证明热路径调用某函数，需要更强证据：

\begin{itemize}
  \item 调用点反汇编中存在 \instruction{call} 到该符号或其 PLT 入口。
  \item 运行时采样显示指令指针落在该函数内。
  \item GDB 断点在该函数上能被当前输入触发。
  \item 计数器或 tracing 记录该函数调用次数。
\end{itemize}

证据强度不同，结论也不同。 \code{nm} 只能支持“函数存在”；反汇编调用点支持“机器码包含调用”；GDB 或采样支持“这次运行走到了那里”。报告必须匹配证据强度。

\section{从反汇编恢复函数边界的谨慎规则}

反汇编工具通常根据符号表、section、调试信息和指令流显示函数。但函数边界并不是天然写在每条机器码旁边。优化、链接器布局、热冷拆分、尾调用、异常表和手写汇编都可能让边界看起来不直观。

恢复函数边界时，先看符号表。全局函数、局部函数、弱符号、PLT 条目、thunk 和冷路径符号可能都出现在输出中。再看 \code{.type} 和 \code{.size} 信息，判断工具认为函数范围从哪里到哪里。若没有大小信息，反汇编器可能用下一个符号作为边界，这在手写汇编或 stripped 二进制中可能不可靠。

其次看控制流。函数入口通常是调用目标或导出符号；函数结束可能是 \instruction{ret}，也可能是尾调用 \instruction{jmp}，也可能是不返回调用。一个 \instruction{ret} 后面出现字节，不代表同一函数继续执行；可能是下一个函数、对齐填充、跳转表或冷路径。反过来，一个函数也可能有多个返回点。

\subsection{尾调用让调用栈外观改变}

如果函数最后只返回另一个函数的结果，优化器可能把 \instruction{call} 加 \instruction{ret} 改成 \instruction{jmp}。这会让当前函数不再建立普通返回帧。反汇编中你会看到：

\begin{lstlisting}
wrapper:
    jmp target
\end{lstlisting}

这不是缺少 \instruction{ret}，而是尾调用。分析时要问：参数是否已经按目标函数 ABI 放好？当前函数是否已经恢复栈和 callee-saved 寄存器？如果都满足，尾跳合法。调试和 profiling 中，\code{wrapper} 可能不显示为普通调用帧，这也是机器码形态导致的正常现象。

\subsection{热冷拆分让一个函数不连续}

优化器和链接器可能把不常执行的错误路径、异常路径或冷分支放到远离热路径的位置。这称为热冷拆分或类似布局优化。源码上一个函数，在二进制中可能显示为热部分和冷部分，符号名也可能带有 \code{cold}。性能分析时，热路径和冷路径要分开看。冷路径代码大，不代表核心循环大；热路径短，也不代表函数语义简单。

若报告分析函数成本，应说明分析的是热路径、完整函数，还是某个输入触发的路径。把冷错误处理路径和热循环混在一起，会误判代码尺寸、分支结构和指令计数。

\section{汇编阅读中的证据等级}

为了避免过度推断，可以把汇编阅读结论分成四个等级。

第一等级是文本事实。比如“反汇编中存在 \instruction{mov eax, edi}”“目标文件有 \code{R_X86_64_PLT32} 重定位”“符号表中 \code{foo} 是 undefined”。这些可以直接从工具输出看到。

第二等级是机器语义。比如“该指令把 \register{edi} 的低 32 位复制到 \register{eax} 并清零高 32 位”“该 \instruction{call} 目标需要链接器修补”“该内存操作读取 4 字节”。这些来自 ISA 和目标文件规则。

第三等级是源码映射。比如“若 \register{rdi} 是参数 \code{p}，则该 load 读取 \code{p[i]}”“这个 \code{+8} 对应结构体字段 \code{count}”。这些需要函数签名、类型布局和源码上下文。

第四等级是性能解释。比如“这里可能受内存带宽限制”“这个 \instruction{cmov} 可能避免分支错误预测但增加数据依赖”。这些需要测量或至少明确假设。不能把第四等级写成第一等级。

\begin{center}
\begin{tabular}{p{0.18\linewidth}p{0.34\linewidth}p{0.32\linewidth}}
\toprule
等级 & 例子 & 需要的证据 \\
\midrule
文本事实 & 存在某条指令 & 反汇编输出 \\
机器语义 & 指令读写哪些状态 & ISA 语义 \\
源码映射 & 地址对应哪个字段 & 源码类型和布局 \\
性能解释 & 为什么快或慢 & 测量、计数器、对照实验 \\
\bottomrule
\end{tabular}
\end{center}

这个分级能保护报告。比如你可以确定地写“这里有一次 4 字节 load”，谨慎地写“结合源码它应是读取 \code{count} 字段”，再把“它可能因随机访问造成 cache miss”列为待验证假设。

\section{读优化汇编时的十个固定问题}

每次分析优化构建，建议按下面十个问题检查：

\begin{enumerate}
  \item 这个函数是否仍有独立符号，还是被内联或删除。
  \item 调用点是否真的调用它，还是调用被优化掉。
  \item 入口寄存器是否能从 ABI 和函数签名解释。
  \item 局部变量是否仍有对应存储，还是只在寄存器或完全消失。
  \item 循环是否仍存在，是否被展开、向量化或常量折叠。
  \item 分支是否仍是分支，还是变成 \instruction{setcc}/\instruction{cmovcc}/mask。
  \item 内存访问是否对应源码对象，是否有额外 reload 或 store。
  \item 是否存在外部调用、PLT/GOT、间接调用或尾调用。
  \item 反汇编来自目标文件、可执行文件还是运行时映射。
  \item 当前结论是事实、源码映射还是性能假设。
\end{enumerate}

这十个问题看似繁琐，但它们能防止最常见的误读：把源码函数当成机器函数，把符号存在当成执行路径，把 \code{-O0} 形态当成性能路径，把一次地址公式当成唯一源码写法，把反汇编片段当成完整性能证据。

\section{综合案例：读一个陌生函数的完整过程}

假设你只拿到下面反汇编和函数声明：

\begin{lstlisting}[language=C++]
struct Rec {
    std::int32_t id;
    std::int32_t score;
    std::int64_t total;
};

extern "C" std::int64_t update_score(Rec* p, int i, int delta);
\end{lstlisting}

反汇编核心片段：

\begin{lstlisting}
movsxd rax, esi
shl    rax, 4
add    dword ptr [rdi + rax + 4], edx
movsxd rcx, dword ptr [rdi + rax + 4]
add    qword ptr [rdi + rax + 8], rcx
mov    rax, qword ptr [rdi + rax + 8]
ret
\end{lstlisting}

第一轮写机器事实。入口 \register{rdi} 是第一个参数，\register{esi} 是第二个参数低 32 位，\register{edx} 是第三个参数。第一条把 \code{i} 符号扩展到 64 位，第二条把它乘以 16。第三条对地址 \code{rdi + i*16 + 4} 做 4 字节 read-modify-write，加上 \code{delta}。第四条从同一地址读取 4 字节并符号扩展到 64 位。第五条把该值加到地址 \code{rdi + i*16 + 8} 的 8 字节整数。第六条读取该 8 字节整数作为返回值。

第二轮结合布局。 \code{Rec} 的字段 offset 很可能是：

\begin{lstlisting}
id    offset 0,  size 4
score offset 4,  size 4
total offset 8,  size 8
sizeof(Rec) = 16
\end{lstlisting}

因此 \code{rdi + i*16 + 4} 对应 \code{p[i].score}，\code{rdi + i*16 + 8} 对应 \code{p[i].total}。注意这个结论依赖结构体布局；若源码类型不同，同样地址公式可能有别的解释。

第三轮写 C-like 语义：

\begin{lstlisting}[language=C++]
p[i].score += delta;
p[i].total += p[i].score;
return p[i].total;
\end{lstlisting}

第四轮写边界。汇编不能证明 \code{i} 在范围内，不能证明 \code{p} 非空，不能证明 \code{p} 指向至少 \code{i+1} 个对象，也不能证明并发访问安全。它只能证明当前机器码按该地址公式访问内存。C++ 合法性来自调用者前置条件。

第五轮写性能假设。该函数对同一结构体元素做一次 4 字节 RMW、一次 4 字节 load、一次 8 字节 RMW、一次 8 字节 load。若在循环中随机访问 \code{p[i]}，可能受 cache 影响；若顺序访问，结构体大小 16 对 cache line 比较规整。但这只是静态假设，不能替代测量。

这个案例展示了完整报告的层次：机器事实、布局解释、伪代码、合法性边界、性能假设。缺任何一层，报告都不完整。

\section{自测清单：汇编语法是否真正过关}

本章收束时，应能独立回答：

\begin{enumerate}
  \item \filepath{.s} 文件和 \code{objdump} 输出分别来自哪个阶段。
  \item 伪指令、标签和 CPU 指令如何区分。
  \item \code{.text}、\code{.data}、\code{.rodata}、\code{.bss} 的作用。
  \item Intel 语法中目的操作数和源操作数顺序。
  \item 内存操作数的 base、index、scale、displacement 如何组合。
  \item 写 32 位寄存器为什么会影响 64 位寄存器。
  \item \instruction{lea} 为什么不一定访问内存。
  \item flags 生产者和消费者如何配对。
  \item 函数入口寄存器如何从 ABI 推导，而不是从第一条指令猜。
  \item 哪些结论能从反汇编直接看出，哪些必须依赖源码类型。
\end{enumerate}

\section{深入讲解：反汇编不是反编译}

反汇编把机器码字节翻译成人类可读的指令文本；反编译试图恢复高级语言结构。两者目标不同。初学者常把反汇编当成“低级版源码”，于是期待每条指令都能对应一行 C++。这种期待会导致误读。

机器码没有 C++ 的变量名、作用域、模板、重载、对象生命周期、异常语义和完整类型信息。调试信息可以补充一部分，但优化后仍然可能不完整。反汇编能告诉你寄存器、内存、跳转和调用；不能单独告诉你某个地址一定属于哪个 C++ 对象，也不能证明某个源码变量仍然存在。

因此，反汇编阅读应保持双向克制。向下，要尊重机器事实：指令读写什么、地址怎么算、宽度是多少。向上，要谨慎恢复源码：只有结合函数签名、类型布局、符号、调试信息和上下文，才能写 C-like 伪代码。伪代码是分析产物，不是证据本身。

\subsection{为什么工具显示会不同}

不同反汇编工具可能显示不同文本。同一条机器码，\code{objdump} 和 \code{llvm-objdump} 的符号命名、立即数格式、地址显示、重定位旁注、伪指令风格可能不同。Intel 和 AT\&T 语法也不同。工具显示不同不代表机器码不同。

严肃比较时，要保存机器码字节和工具命令。若两个输出指令文本不同，先比较字节；若字节相同，只是显示差异。若字节不同，再回到构建命令、链接阶段和工具版本查原因。不要只凭视觉差异判断程序变化。

\section{深入讲解：汇编学习要避免“指令表崇拜”}

很多资料把汇编学习变成背指令表。指令表有用，但不是核心。真正的核心是状态变化和上下文。同一条 \instruction{mov}，在函数入口可能是参数搬运，在循环中可能是 load，在返回前可能是准备返回值，在 ABI 边界可能是保存调用者状态。同一条 \instruction{lea}，可能是地址计算，也可能是整数乘加。

学习一条指令时，要同时学：

\begin{itemize}
  \item 它读哪些状态。
  \item 它写哪些状态。
  \item 它是否访问内存。
  \item 它是否改 flags。
  \item 它的操作数宽度如何决定。
  \item 它常出现在什么编译器模式里。
  \item 它在 ABI 或控制流上下文中可能承担什么角色。
\end{itemize}

这样学习，指令才会进入工程能力。单独背“\instruction{mov} 是传送指令”“\instruction{lea} 是装入有效地址”，反而容易形成误解。

\section{读汇编要先建立坐标系}

很多人第一次打开反汇编，会从第一条指令开始逐行翻译。这样很容易迷路，因为你还不知道这段文本处在什么坐标系里。读汇编前要先建立四个坐标：文件坐标、函数坐标、控制流坐标和数据坐标。

文件坐标回答“这段指令来自哪个二进制”。同一个源码可以生成多个目标文件和多个可执行文件，目标文件中的地址、重定位和符号与最终可执行文件不同。报告中应写清工具命令、输入文件、构建参数和是否经过链接。只贴一段反汇编，不说明来源，读者无法判断重定位、地址和符号是否可信。

函数坐标回答“函数边界在哪里”。符号表、标签、\code{.type}、\code{.size}、反汇编分隔和调用目标都可以帮助确认函数范围。但优化后函数可能被内联、拆分、合并或冷路径外提。若函数标签不存在，不一定表示源码函数没有执行；若存在标签，也不一定表示它没有被其他路径绕过。函数坐标要结合构建和优化级别判断。

控制流坐标回答“执行路径如何移动”。线性地址顺序只是布局，不是语义顺序。条件跳转、无条件跳转、fallthrough、call、ret、跳转表和异常边都会改变路径。应先标出基本块，再读块内指令。否则容易把冷路径当成热路径，或者把跳过的代码当成必经代码。

数据坐标回答“每个寄存器和内存地址扮演什么角色”。入口寄存器来自 ABI，栈槽可能是局部变量或保存寄存器，全局地址可能通过 RIP-relative 访问，结构体字段通过 displacement 暴露。数据坐标不是一次性确定的，寄存器角色会随基本块变化。读汇编报告应记录角色变化，而不是给某个寄存器贴永久标签。

\section{综合案例：从陌生函数恢复最小语义}

假设你拿到一个没有源码上下文的小函数，只知道它属于 x86-64 System V 程序。高质量阅读可以按六步进行。第一步，确认入口和出口。查看函数标签、\instruction{ret}、尾调用 \instruction{jmp}、栈调整和可能的不返回调用。第二步，确认 ABI 入口寄存器。若看到 \register{rdi}、\register{rsi}、\register{rdx} 在函数开始被使用，先把它们标成候选参数，而不是直接命名成具体变量。

第三步，分基本块。每个条件跳转和目标标签都可能产生新块。把块画出来后，先写路径条件，例如“若 \register{esi} 小于等于零则返回零”。第四步，恢复数据流。找出循环变量、累加器、指针递增、load/store 宽度和地址公式。第五步，结合类型猜测写 C-like 伪代码，但要标注不确定性。比如“可能是 \code{int const*}，因为 load 宽度为 4 且下标步长为 4”。

第六步，写边界。反汇编不能证明指针非空，不能证明长度合法，不能证明别名关系，也不能证明源语言类型。它只能证明这份机器码按某些地址公式和控制路径执行。若报告把伪代码写得像源码原文，就过度解释了。正确写法是“在假设第一个参数是指向 32 位整数的指针、第二个参数是元素个数时，该函数等价于求和非负长度数组”。

\subsection{把不确定性写出来}

底层报告不是越肯定越好，而是证据和结论匹配越好越专业。你可以把不确定性分级。一级证据是机器事实：某指令读 \register{rdi} 加偏移，某块由 \instruction{jle} 跳入。二级解释是 ABI 和布局推断：\register{rdi} 可能是第一个指针参数，偏移 8 可能是结构体第二个字段。三级假设是源码恢复：这个字段可能叫 \code{size} 或 \code{total}。不同层级不能混写。

当证据不足时，可以列出需要补的观察：函数声明、调用点、调试信息、结构体定义、更多输入输出样本、GDB 单步、反编译辅助、符号表或重定位信息。写清需要什么证据，说明你知道当前结论的边界。

\section{汇编阅读证据结构}

每次汇编阅读都应固定保留以下证据：

\begin{enumerate}
  \item 目标身份：源码提交、构建命令、二进制路径、工具命令。
  \item 函数身份：符号名、地址范围、是否内联、是否外部可见。
  \item ABI 假设：平台、调用约定、入口参数、返回值位置。
  \item 基本块：入口、分支、循环、退出、冷路径。
  \item 寄存器角色：每个阶段主要寄存器保存什么。
  \item 内存访问：宽度、地址公式、可能对象、对齐和边界前提。
  \item 控制依赖：哪些检查保护哪些危险操作。
  \item C-like 伪代码：只写证据支持的最小语义。
  \item 性能假设：指令类型、依赖链、分支、load/store 和可能瓶颈。
  \item 限制：哪些结论不能从当前反汇编证明。
\end{enumerate}

固定结构不是为了增加形式感，而是避免遗漏关键问题。很多错误汇编解读只写了伪代码，没有写 ABI；只写了指令含义，没有写控制流；只写了性能猜测，没有写输入和边界。证据结构把这些层次分开，结论才不会越过反汇编能够证明的范围。

\section{常见视觉陷阱}

反汇编有一些视觉陷阱。第一，地址相邻不代表路径相邻。编译器可能把冷路径放在热路径之后，也可能把多个函数的冷块集中到其他 section。第二，短指令不一定便宜，长指令不一定慢。性能要看依赖、吞吐、内存和微架构。第三，\instruction{lea} 看起来像地址指令，但常用于整数乘加。第四，\instruction{xor eax, eax} 不是逻辑意图上的“异或”，常是清零。

第五，栈槽不一定对应源码局部变量。优化后栈槽可能用于溢出寄存器、保存调用参数、对齐、异常处理或临时存储。第六，\code{nop} 不一定是无意义垃圾，可能用于对齐或热补丁。第七，同一个源码变量可能在不同路径上位于不同寄存器，甚至不存在。第八，反汇编中的符号名可能来自最近的符号，不一定是精确标签。

这些陷阱的共同原因是：反汇编是机器码视角，不是源码视角。分析时要尊重它的事实，同时控制向源码层推断的欲望。

\section{符号、重定位和地址注释}

汇编阅读不能只看指令助记符，还要看符号和重定位。目标文件中的许多地址还没有最终确定，反汇编旁边的重定位记录说明链接器以后会填什么。若一条指令访问 \code{rip + 0}，旁边有指向某个全局符号的重定位，不能按字面把偏移零解释成真实地址。要结合 \code{objdump -dr} 的重定位注释。

函数调用也是如此。目标文件里对外部函数的 \instruction{call} 可能显示相对偏移占位，重定位记录指向符号；最终可执行文件中可能变成 PLT 调用，也可能在静态链接或 LTO 后变成直接地址或被内联。阅读目标文件和阅读最终可执行文件，结论层级不同。报告中要说明你看的是哪一层。

RIP-relative 寻址是 x86-64 常见模式。它让代码可以位置无关地访问附近数据或 GOT 表项。看到 \code{[rip + disp]}，要计算的是下一条指令地址加 displacement，而不是当前指令起始地址加 displacement。初学者常在这里算错地址。若有符号注释，以符号和重定位为准；若没有，就用地址手工验证。

\section{调试信息和反汇编的配合}

带 \code{-g} 的二进制可以把指令映射回源码行，但这种映射不是完美逆变换。优化后，一条源码行可能对应多段指令，一条指令可能服务多个源码表达式，变量可能被优化掉或位于多个位置。GDB 显示的源码行是线索，不是完整解释。

阅读时可以先用调试信息定位函数和源码区域，再切换到指令级视角。若 GDB 显示变量为 optimized out，不要立刻认为变量不存在；它可能被常量传播、拆成寄存器片段，或者根本不需要存储。若源码单步跳来跳去，可能是优化重排和内联造成的。此时应回到反汇编和寄存器状态。

调试信息还可能和二进制不匹配。若符号文件来自另一构建，源码行会误导。正式报告应记录 build id 或二进制校验和，确保调试信息和机器码对应。反汇编阅读的第一原则仍然是：先确认对象身份。

\section{从汇编到问题单}

真实工程中，汇编阅读常常服务一个问题单：为什么崩溃，为什么慢，为什么调用错库，为什么 ABI 破坏，为什么某个函数消失。报告应围绕问题，不必把每条指令都翻译。先列与问题相关的块，再说明它们如何支持或推翻假设。

例如性能问题关注热循环、load/store、分支和调用；崩溃问题关注当前指令、寄存器、内存地址和保护条件；ABI 问题关注入口寄存器、栈对齐和保存寄存器；链接问题关注符号和重定位。目的不同，阅读重点不同。会读汇编，不是把所有指令翻成中文，而是从机器事实中提取和问题相关的证据。

\section{汇编注释应该解释层次，不要复述指令}

给汇编写注释时，低质量注释只是把指令翻译成中文，例如“把 \register{eax} 赋值为 0”。这种注释帮助很小。高质量注释应解释层次：这个寄存器此刻代表哪个逻辑值，这个分支保护什么条件，这个地址公式对应哪个字段，这个调用跨过哪个 ABI 边界，这个块为什么是冷路径。

例如：

\begin{lstlisting}
xor eax, eax        ; return value = 0 for n <= 0 path
test esi, esi       ; check n before any load from p
jle .Lreturn        ; zero/negative length does not access memory
\end{lstlisting}

这里注释说明了控制流保护关系，而不只是复述 \instruction{test} 和 \instruction{jle}。注释应说明为什么这几条指令重要。汇编报告中的注释也应如此：解释证据如何连接到语义和边界。

\section{从手写汇编反推编译器生成汇编}

学习汇编有两条路线：读编译器生成的汇编，写少量手写汇编。手写能帮助理解语法、标签、section、符号和 ABI；读编译器汇编能帮助理解优化、寄存器分配、控制流和真实代码形态。两者要互相校准。只写手写汇编，容易形成不符合编译器输出的风格；只读编译器汇编，容易缺少对目标文件结构的掌控。

手写最小函数后，可以马上写等价 C++，比较两者目标文件。看符号、section、\code{.note.GNU-stack}、\code{.type}、\code{.size}、栈帧和返回值是否一致。这个练习能把“汇编文本”与“可链接目标文件”联系起来，为第 13 和第 14 章做准备。

\section{阅读优化汇编的耐心}

优化汇编常常比 \code{-O0} 难读，因为源码变量消失、循环形态改变、分支反转、常量折叠、内联和公共子表达式消除都会发生。不要把这种变化理解成“编译器乱写”。优化器是在语言合同允许范围内重写程序，使机器执行更合适。

阅读时可以先用 \code{-O0} 建立源码到机器的粗映射，再用 \code{-O2} 或 \code{-O3} 看真实发布形态。两者都要看，但回答不同问题。\code{-O0} 帮助学习语义，\code{-O3} 帮助分析性能。若报告要支持性能结论，最终必须回到优化构建；若报告要解释语言结构，可以用低优化辅助理解。

\section{综合复盘：从一段汇编写出三层结论}

可以拿任意一个短函数练习三层结论。第一层是机器事实：指令、寄存器、内存宽度、地址公式和跳转。第二层是源码解释：在某些类型假设下，它可能对应什么 C-like 语义。第三层是工程判断：它是否满足 ABI，是否有边界风险，是否有性能假设。三层结论必须分开写。

例如一条 \code{mov eax, dword ptr [rdi + rsi*4]}，机器事实是从地址读取 4 字节到 \register{eax}；源码解释可能是读取 \code{int} 数组第 \code{i} 个元素；工程判断还要问 \register{rdi} 是否非空、\register{rsi} 是否在范围内、读取是否受检查保护、访问是否顺序。只写“读取数组元素”省略了很多前提。

这个练习强调克制。汇编阅读最重要的品质不是想象力，而是知道哪些结论有证据，哪些只是合理猜测。

\section{汇编阅读的长期方法}

每周选择一个小函数，保存源码、构建命令、反汇编和一页分析。函数可以很简单：加法、数组求和、结构体字段更新、短路判断、switch、函数调用。简单函数重复做，能让模式稳定进入直觉。等模式熟悉后，再读复杂函数就不会被局部指令淹没。

长期积累还要保留失败样例。哪些指令曾被误解，哪些条件码曾读反，哪些地址公式曾算错，哪些函数边界曾判断错，都应该写进错误笔记。汇编能力不是靠一次性背完指令表形成，而是靠不断校正误读形成。

\section{汇编报告的审阅方式}

审阅一份汇编报告时，不要先挑指令翻译是否漂亮，而要先看证据层次。报告是否说明二进制身份，是否说明平台 ABI，是否区分机器事实和源码推断，是否写出无法证明的部分，是否把控制流和数据流连起来。如果这些层次缺失，即使单条指令翻译正确，报告仍然不可靠。

其次看边界。报告是否说明指针和下标的合法性来自哪里，是否说明结构体 offset 的依据，是否说明优化级别和内联对函数边界的影响，是否说明性能假设需要测量。汇编报告最常见的问题不是少翻译一条指令，而是把推断写成事实。审阅时要鼓励作者把不确定性写出来。

最后看可复现性。别人能否用同一命令生成同一段反汇编，能否找到同一函数，能否验证同一地址公式。如果不能，报告就只是一次口头解释。汇编学习从个人能力变成工程能力，关键就在可复现性。

\section{实验：建立第一份汇编阅读报告}

\begin{labbox}
创建 \filepath{labs/ch10_assembly_syntax/basic_reading.cpp}：

\begin{lstlisting}[language=C++]
#include <cstddef>
#include <cstdint>

extern "C" int add2(int x, int y)
{
    return x + y;
}

extern "C" int mix(int x)
{
    return x * 5 + 17;
}

extern "C" std::uint8_t load_u8(std::uint8_t const* p, std::size_t i)
{
    return p[i];
}

extern "C" int load_i32(int const* p, int i)
{
    return p[i + 2];
}
\end{lstlisting}

生成汇编和反汇编：

\begin{lstlisting}[language=bash]
clang++ -std=c++20 -O0 -S -masm=intel basic_reading.cpp -o basic_reading_O0.s
clang++ -std=c++20 -O3 -S -masm=intel basic_reading.cpp -o basic_reading_O3.s
clang++ -std=c++20 -O3 -c basic_reading.cpp -o basic_reading.o
objdump -d -Mintel basic_reading.o > basic_reading.objdump.txt
\end{lstlisting}

交付物：

\begin{enumerate}
  \item 对每个函数列出 ABI 参数寄存器和返回值寄存器。
  \item 从 \code{basic_reading_O3.s} 中摘出每个函数的核心指令。
  \item 对每条核心指令写“读什么、写什么、是否访问内存、是否更新 flags”。
  \item 对 \code{load_u8} 解释为什么可能出现 \instruction{movzx}。
  \item 对 \code{mix} 解释编译器是否使用 \instruction{lea}。
  \item 对 \code{load_i32} 写出地址公式，并说明 \code{+ 8} 从哪里来。
  \item 从 \code{objdump} 中记录机器码字节，说明汇编文本和机器码字节的关系。
\end{enumerate}
\end{labbox}

\section{实验：Intel 和 AT\&T 语法互译}

\begin{labbox}
用同一个源码生成两种语法：

\begin{lstlisting}[language=bash]
clang++ -std=c++20 -O3 -S -masm=intel basic_reading.cpp -o intel.s
clang++ -std=c++20 -O3 -S basic_reading.cpp -o att.s
\end{lstlisting}

交付物：

\begin{enumerate}
  \item 从两个文件中各选 8 条对应指令。
  \item 写成三列表：Intel、AT\&T、语义。
  \item 标出操作数顺序差异。
  \item 标出内存寻址写法差异。
  \item 标出宽度写法差异，例如 \code{dword ptr} 和 \code{movl}。
\end{enumerate}

报告中必须包含下面这种格式：

\begin{center}
\begin{tabular}{p{0.29\linewidth}p{0.29\linewidth}p{0.28\linewidth}}
\toprule
Intel & AT\&T & 语义 \\
\midrule
\code{mov eax, edi} & \code{movl \%edi, \%eax} & \code{eax = edi} \\
\code{add eax, esi} & \code{addl \%esi, \%eax} & \code{eax += esi} \\
\bottomrule
\end{tabular}
\end{center}
\end{labbox}

\section{实验：\code{-O0} 和 \code{-O3} 对比}

\begin{labbox}
创建 \filepath{labs/ch10_assembly_syntax/o0_o3_compare.cpp}：

\begin{lstlisting}[language=C++]
extern "C" int formula(int x)
{
    int a = x + 1;
    int b = a * 3;
    int c = b - x;
    return c + 7;
}
\end{lstlisting}

生成：

\begin{lstlisting}[language=bash]
clang++ -std=c++20 -O0 -S -masm=intel o0_o3_compare.cpp -o formula_O0.s
clang++ -std=c++20 -O3 -S -masm=intel o0_o3_compare.cpp -o formula_O3.s
\end{lstlisting}

交付物：

\begin{enumerate}
  \item 注释 \code{-O0} 中的栈帧建立、局部变量栈槽和返回路径。
  \item 注释 \code{-O3} 中的代数化简结果。
  \item 用数学式证明 \code{formula(x)} 可以化简成什么。
  \item 解释为什么不能用 \code{-O0} 汇编评估性能。
  \item 说明 \code{-O0} 对调试有什么价值。
\end{enumerate}
\end{labbox}

\section{实验：符号、section 和重定位证据}

\begin{labbox}
目标：把 \filepath{.s}、\filepath{.o}、符号表和重定位表连起来看，理解目标文件里哪些地址已经确定，哪些还要等链接器修补。

创建 \filepath{labs/ch10_assembly_syntax/relocation_a.cpp}：

\begin{lstlisting}[language=C++]
#include <cstdio>

extern "C" int external_add(int, int);

int global_counter = 7;

extern "C" int call_external(int x)
{
    return external_add(x, global_counter);
}

extern "C" void print_value(int x)
{
    std::printf("value=%d\n", x);
}
\end{lstlisting}

创建 \filepath{labs/ch10_assembly_syntax/relocation_b.cpp}：

\begin{lstlisting}[language=C++]
extern "C" int external_add(int a, int b)
{
    return a + b;
}
\end{lstlisting}

先只编译第一个目标文件：

\begin{lstlisting}[language=bash]
clang++ -std=c++20 -O2 -fPIC -S -masm=intel relocation_a.cpp -o relocation_a.s
clang++ -std=c++20 -O2 -fPIC -c relocation_a.cpp -o relocation_a.o
objdump -drwC -Mintel relocation_a.o > relocation_a.objdump.txt
readelf -S -s -r relocation_a.o > relocation_a.readelf.txt
nm -C relocation_a.o > relocation_a.nm.txt
\end{lstlisting}

再把两个目标文件链接成可执行文件：

\begin{lstlisting}[language=bash]
clang++ -std=c++20 -O2 relocation_a.o relocation_b.cpp -o relocation_app
objdump -drwC -Mintel relocation_app > relocation_app.objdump.txt
readelf -S -s -r relocation_app > relocation_app.readelf.txt
nm -C relocation_app > relocation_app.nm.txt
\end{lstlisting}

交付物：

\begin{enumerate}
  \item 在 \code{relocation_a.s} 中标出 \code{.text}、\code{.data}、\code{.rodata} 或相关 section 切换。
  \item 用 \code{nm} 判断 \code{external_add}、\code{printf}、\code{global_counter} 哪些是已定义符号，哪些是未定义符号。
  \item 在 \code{objdump -drwC} 输出中找出访问 \code{global_counter} 的指令及其重定位旁注。
  \item 找出调用 \code{external_add} 或 \code{printf} 的 \instruction{call}，解释目标文件阶段为什么可能显示占位位移。
  \item 用 \code{readelf -r} 把重定位 offset、重定位类型、符号名和反汇编中的指令位置对应起来。
  \item 对比目标文件和链接后可执行文件，说明哪些地址或调用目标发生了变化。
  \item 说明 \code{-fPIC} 对全局变量访问形态可能产生什么影响。
\end{enumerate}

报告要求：不能只贴输出。每个结论都要写清“证据来自哪个文件、哪个命令、哪一行输出”。例如“\code{external_add} 在 \code{relocation_a.o} 中是 undefined symbol，证据是 \code{nm -C relocation_a.o} 显示 \code{U external_add}；对应 \instruction{call} 的修补证据来自 \code{objdump -drwC} 中的 relocation 行”。
\end{labbox}

\section{实验：手写最小汇编目标文件}

\begin{labbox}
目标：亲手写一个最小 \filepath{.S} 文件，观察标签、全局符号、函数类型、大小和不可执行栈标记如何进入目标文件。

创建 \filepath{labs/ch10_assembly_syntax/handwritten.S}：

\begin{lstlisting}
.intel_syntax noprefix
.text
.globl asm_add2
.type asm_add2,@function
asm_add2:
    mov eax, edi
    add eax, esi
    ret
.size asm_add2, .-asm_add2

.section .note.GNU-stack,"",@progbits
\end{lstlisting}

创建 \filepath{labs/ch10_assembly_syntax/use_handwritten.cpp}：

\begin{lstlisting}[language=C++]
#include <cstdio>

extern "C" int asm_add2(int, int);

int main()
{
    std::printf("%d\n", asm_add2(20, 22));
}
\end{lstlisting}

构建并观察：

\begin{lstlisting}[language=bash]
clang++ -c handwritten.S -o handwritten.o
clang++ -std=c++20 -c use_handwritten.cpp -o use_handwritten.o
clang++ handwritten.o use_handwritten.o -o use_handwritten
./use_handwritten
objdump -drwC -Mintel handwritten.o > handwritten.objdump.txt
readelf -S -s handwritten.o > handwritten.readelf.txt
nm -C handwritten.o > handwritten.nm.txt
\end{lstlisting}

交付物：

\begin{enumerate}
  \item 解释 \code{.intel_syntax noprefix} 改变了什么。
  \item 解释 \code{.globl asm_add2} 为什么对跨文件链接必要。
  \item 解释 \code{.type} 和 \code{.size} 对工具输出有什么帮助。
  \item 在 \code{objdump} 中记录 \instruction{mov}、\instruction{add}、\instruction{ret} 的机器码字节。
  \item 在 \code{readelf -s} 中找到 \code{asm_add2} 的符号类型、section 和大小。
  \item 说明 \code{.note.GNU-stack} 的目的，以及为什么它不是 CPU 执行路径。
  \item 删除 \code{.globl asm_add2} 后重新构建，记录链接错误，再恢复该行。
\end{enumerate}

这个实验故意只写一个加法函数。目标不是展示复杂汇编，而是让你看到：可链接的汇编函数不仅有三条机器指令，还需要正确的符号可见性和目标文件元信息。
\end{labbox}

\section{本章作业}

\begin{exercisebox}
\subsection*{作业 1：逐条注释 20 条汇编}

从本章实验生成的 \code{-O3} 汇编中选择至少 20 条真实指令。每条指令必须写：

\begin{itemize}
  \item 指令文本。
  \item 操作数类型：立即数、寄存器、内存。
  \item 读取的寄存器或内存。
  \item 写入的寄存器或内存。
  \item 访问宽度。
  \item 是否更新 flags。
  \item 你认为它对应的 C++ 语义。
\end{itemize}

\subsection*{作业 2：地址公式恢复}

对下面汇编片段恢复可能的地址公式和 C-like 伪代码：

\begin{lstlisting}
movsxd rax, esi
mov    ecx, dword ptr [rdi + rax*4 + 12]
add    ecx, edx
mov    dword ptr [rdi + rax*4 + 12], ecx
mov    eax, ecx
ret
\end{lstlisting}

要求至少给出两种可能来源：

\begin{enumerate}
  \item \code{int*} 数组访问版本。
  \item 结构体数组字段访问版本。
\end{enumerate}

每种版本都必须解释 \code{+ 12} 的来源。

\subsection*{作业 3：语法互译}

把下面 Intel 语法翻译成 AT\&T 语法，并写出语义：

\begin{lstlisting}
mov rax, rdi
add eax, esi
movzx eax, byte ptr [rdi + rsi]
mov qword ptr [rsp + 8], rax
lea rax, [rdi + rsi*8 + 16]
\end{lstlisting}

再把下面 AT\&T 语法翻译成 Intel 语法：

\begin{lstlisting}
movl %edi, %eax
addl %esi, %eax
movzbl (%rdi,%rsi), %eax
movq %rax, 8(%rsp)
leaq 16(%rdi,%rsi,8), %rax
\end{lstlisting}

\subsection*{作业 4：机器码和指令长度}

用 \code{objdump} 记录至少 10 条指令的机器码字节。回答：

\begin{enumerate}
  \item 哪些指令是 1 字节？
  \item 哪些指令长度超过 4 字节？
  \item 立即数和 displacement 如何让指令变长？
  \item 为什么 x86 前端需要处理变长指令？
\end{enumerate}

\subsection*{作业 5：目标文件证据链}

选择一个包含外部函数调用、字符串字面量、全局变量访问的小程序，分别生成 \filepath{.s}、\filepath{.o} 和可执行文件。提交一份目标文件证据链报告，必须包含：

\begin{itemize}
  \item \code{objdump -drwC -Mintel} 输出中至少 5 条带机器码字节的指令。
  \item \code{readelf -S} 中相关 section 的名称、类型和权限。
  \item \code{readelf -r} 中至少 2 条重定位记录。
  \item \code{nm -C} 中至少 2 个已定义符号和 2 个未定义符号。
  \item 对一个 \instruction{call} 占位位移的解释。
  \item 对一个 \code{[rip + ...]} 全局或常量访问的解释。
  \item 目标文件阶段和链接后阶段地址语境的区别。
\end{itemize}

报告必须明确标注“确定事实、合理推论、待验证假设”。例如“重定位类型是什么”是工具输出事实；“它可能通过 PLT 调用动态库函数”是结合构建方式的推论；“这部分开销影响性能”则需要运行时测量，不能仅凭反汇编断言。

\subsection*{作业 6：小型汇编阅读报告}

选择一个你自己写的 10 到 20 行 C++ 函数，要求包含：

\begin{itemize}
  \item 至少一个数组访问。
  \item 至少一个结构体字段访问。
  \item 至少一个整数表达式。
  \item 至少一个条件判断或比较。
\end{itemize}

生成 \code{-O0} 和 \code{-O3} 汇编。提交不少于 2000 字报告，必须包含：

\begin{itemize}
  \item 源码。
  \item 编译命令。
  \item 函数签名和 ABI 参数寄存器。
  \item \code{-O0} 和 \code{-O3} 汇编对比。
  \item 至少 15 条指令逐条注释。
  \item 所有内存访问的地址公式。
  \item 至少 5 条机器码字节记录。
  \item 你遇到的一个误判，以及如何纠正。
\end{itemize}
\end{exercisebox}

\section{验收标准}

本章收束时，应能做到：

\begin{itemize}
  \item 熟练区分 Intel 和 AT\&T 语法。
  \item 识别立即数、寄存器和内存操作数。
  \item 解释通用寄存器别名和 32 位写入清零高位。
  \item 从内存操作数写出有效地址公式。
  \item 解释 \code{byte ptr}、\code{dword ptr}、\code{qword ptr}、\code{xmmword ptr} 的访问宽度。
  \item 解释 \instruction{lea} 为什么不访问内存。
  \item 区分 \code{.text}、\code{.rodata}、\code{.data}、\code{.bss} 的基本用途。
  \item 说明 \code{.globl}、\code{.type}、\code{.size}、局部标签和全局符号的作用。
  \item 用 \code{objdump -drwC}、\code{readelf -r}、\code{nm -C} 交叉验证外部调用和全局访问。
  \item 解释目标文件阶段的重定位占位为什么不能直接当作最终地址。
  \item 对简单函数生成 \code{-O0} 和 \code{-O3} 汇编，并逐条注释核心指令。
  \item 从反汇编中记录机器码字节，知道汇编文本最终对应指令编码。
  \item 不再把 \code{-O0} 汇编当作性能路径。
\end{itemize}

如果这些能力稳定，本册后面的整数指令、控制流、ABI 和 C++ 汇编互操作才有可靠基础；第二册的自动向量化和 AVX microkernel 也才真正读得进去。
